\SGAchapter{XV}{Frobenius 态射与 L 函数的有理性}{Morphisme de Frobenius et rationalité des fonctions L}
\begin{center}C. Houzel\end{center}
\SourcePage{454}
\SGAsection{1}{Frobenius 态射}
\textbf{1. \Term{frobenius-endomorphism}{Frobenius 自同态}。}
\begin{SGAresult}{定义}{1}
设 $p$ 为素数。若 $p\cdot1_{\mathcal O_X}=0$，就称\Term{prescheme}{预概形} $X$ 的\Term{characteristic}{特征}为 $p$。
\end{SGAresult}
这一条件\Term{equivalence}{等价}于：唯一的\Term{morphism}{态射} $X\to\operatorname{Spec}(\mathbb Z)$ 经过 $\operatorname{Spec}(\mathbb F_p)\to\operatorname{Spec}(\mathbb Z)$ \Term{resolution}{分解}；还应指出，这种\Term{resolution}{分解}是唯一的。

以 $(\mathrm{Sch})_p$ 表示\Term{prescheme}{预概形}\Term{category}{范畴}中由\Term{characteristic}{特征}为 $p$ 的\Term{prescheme}{预概形}组成的\Term{full-subcategory}{满子范畴}。根据上述说明，$(\mathrm{Sch})_p$ 自然\Term{isomorphism}{同构}于 $\mathbb F_p$-\Term{prescheme}{预概形}的\Term{category}{范畴}。

若 $X$ 是\Term{characteristic}{特征}为 $p$ 的\Term{prescheme}{预概形}，则二元组 $\mathrm{fr}_X=(\operatorname{id}_{|X|},\varphi)$ 是 $X$ 的一个\Term{endomorphism}{自同态}，其中 $\varphi:\mathcal O_X\to\mathcal O_X$ 对每个\Term{open-subset}{开集}上的\Term{section}{截面} $f$ 定义为 $\varphi(f)=f^p$；因为 $\varphi$ 显然是环\Term{sheaf}{层} $\mathcal O_X$ 的\Term{endomorphism}{自同态}。
\SourcePage{455}
\begin{SGAresult}{定义}{2}
设 $X$ 是\Term{characteristic}{特征}为 $p$ 的\Term{prescheme}{预概形}。称上述\Term{endomorphism}{自同态} $\mathrm{fr}_X$ 为 $X$ 的 \Term{frobenius-endomorphism}{Frobenius 自同态}。
\end{SGAresult}
\Term{frobenius-endomorphism}{Frobenius 自同态} $\mathrm{fr}_X$ 对\Term{characteristic}{特征}为 $p$ 的\Term{prescheme}{预概形} $X$ 具有\Term{functoriality}{函子性}：对任意此类\Term{prescheme}{预概形}之间的\Term{morphism}{态射} $g:X\to Y$，下图交换：
\[
\begin{tikzcd}
X\arrow[r,"\mathrm{fr}_X"]\arrow[d,"g"']&X\arrow[d,"g"]\\
Y\arrow[r,"\mathrm{fr}_Y"']&Y.
\end{tikzcd}
\]
换言之，我们定义了恒等\Term{functor}{函子} $\operatorname{id}_{(\mathrm{Sch})_p}$ 的\Term{endomorphism}{自同态} $\mathrm{fr}$。

\textbf{2. \Term{relative-frobenius}{相对 Frobenius 态射}。}

设 $S$ 是\Term{characteristic}{特征}为 $p$ 的\Term{prescheme}{预概形}，则所有 $S$-\Term{prescheme}{预概形}也都具有\Term{characteristic}{特征} $p$。若 $X$ 是 $S$-\Term{prescheme}{预概形}，由下图可见，$\mathrm{fr}_X$ 经过\Term{projection}{投影}
\[
\pi_{X/S}:X\times_S(S,\mathrm{fr}_S)\longrightarrow X
\]
\Term{resolution}{分解}，这里 $(S,\mathrm{fr}_S)$ 表示借助\Term{morphism}{态射} $\mathrm{fr}_S$ 视为 $S$-\Term{prescheme}{预概形}的 $S$。将\Term{fiber-product}{纤维积} $X\times_S(S,\mathrm{fr}_S)$ 记作 $X^{(p/S)}$，在不会混淆时也记作 $X^{(p)}$。从 $X$ 到 $X^{(p/S)}$、与 $\pi_{X/S}$ 复合后等于 $\mathrm{fr}_X$ 的唯一 $S$-\Term{morphism}{态射}，记作 $\mathrm{Fr}_{X/S}$。于是得到如下\Term{commutative-diagram}{交换图}，其中方块是\Term{cartesian-square}{笛卡尔方块}：
\SourcePage{456}
\[
\begin{tikzcd}[column sep=large]
X\arrow[d,"g"']&X^{(p)}\arrow[l,"\pi_{X/S}"']\arrow[d,"g^{(p)}"]&X\arrow[l,"\mathrm{Fr}_{X/S}"']\arrow[ll,bend right=27,"\mathrm{fr}_X"']\arrow[dl,"g"]\\
S&S\arrow[l,"\mathrm{fr}_S"']&
\end{tikzcd}
\]
显然，$X^{(p/S)}$ 对 $S$-\Term{prescheme}{预概形} $X$ 具有\Term{functoriality}{函子性}；\Term{functor}{函子} $P_S:X\mapsto X^{(p/S)}$ 就是沿 $\mathrm{fr}_S:S\to S$ 的\Term{base-change}{基变换}。和所有\Term{base-change}{基变换}一样，它与\Term{projective-limit}{投射极限}及和交换，并保持\Term{finite-type}{有限型}、\Term{separated}{分离}、\Term{affine}{仿射}、\Term{projective}{射影}、\Term{flat}{平坦}、\Term{etale}{平展}等性质。它与\Term{projective-limit}{投射极限}交换这一事实表明：若 $X$ 带有某种\Term{algebraic-structure}{代数结构}（例如群、\Term{algebra}{代数}等），则 $X^{(p/S)}$ 自然带有同类的\Term{algebraic-structure}{代数结构}。

利用 $\mathrm{fr}_X$ 对 $X$ 的\Term{functoriality}{函子性}，可见 $\mathrm{Fr}_{X/S}$ 也对 $S$-\Term{prescheme}{预概形} $X$ 具有\Term{functoriality}{函子性}，从而定义了一个自然\Term{morphism}{态射}
\[
\mathrm{Fr}_{X/S}:\operatorname{id}_{(\mathrm{Sch})/S}\longrightarrow P_S.
\]
\begin{SGAresult}{定义}{3}
设 $S$ 是\Term{characteristic}{特征}为 $p$ 的\Term{prescheme}{预概形}，$X$ 是 $S$-\Term{prescheme}{预概形}。称 $S$-\Term{morphism}{态射} $\mathrm{Fr}_{X/S}:X\to X^{(p/S)}$ 为 $X$ 相对于 $S$ 的 \Term{relative-frobenius}{Frobenius 态射}。
\end{SGAresult}
\Term{relative-frobenius}{Frobenius 态射}最初的\Term{form}{形式}性质列于下一命题。
\SourcePage{457}
\begin{SGAresult}{命题}{1}
a)（\Term{transitivity}{传递性}）设 $S$ 是\Term{characteristic}{特征}为 $p$ 的\Term{prescheme}{预概形}，$g:X\to Y$ 是 $S$-\Term{prescheme}{预概形}之间的 $S$-\Term{morphism}{态射}。下图交换，且其中的三个方块都是\Term{cartesian-square}{笛卡尔方块}：
\[
\begin{tikzcd}[column sep=large,row sep=large]
X\arrow[d,"g"']&X^{(p/S)}\arrow[l,"\pi_{X/S}"']\arrow[d,"g^{(p/S)}"']&X^{(p/Y)}\arrow[l,"\pi_{X/Y}"']\arrow[d,"g^{(p/Y)}"]&X\arrow[l,"\mathrm{Fr}_{X/Y}"']\arrow[ll,bend right=25,"\mathrm{Fr}_{X/S}"']\arrow[dl,"g"]\\
Y\arrow[d]&Y^{(p/S)}\arrow[l,"\pi_{Y/S}"']\arrow[d]&Y\arrow[l,"\mathrm{Fr}_{Y/S}"']\arrow[dl]&\\
S&S\arrow[l,"\mathrm{fr}_S"']&&
\end{tikzcd}
\]
b)（\Term{base-change}{基变换}）设 $g:S'\to S$ 是\Term{characteristic}{特征}为 $p$ 的\Term{prescheme}{预概形}之间的\Term{morphism}{态射}，$X$ 是 $S$-\Term{prescheme}{预概形}。存在对 $X$ 具有\Term{functoriality}{函子性}的\Term{canonical}{典范}\Term{isomorphism}{同构}
\[
X^{(p/S)}\times_SS'\simeq(X\times_SS')^{(p/S')}.
\]
若用撇号表示沿 $g$ 的\Term{base-change}{基变换}，即为 $(X^{(p/S)})'\simeq X'^{(p/S')}$；在此\Term{isomorphism}{同构}给出的认同下，还有 $\mathrm{Fr}_{X'/S'}=(\mathrm{Fr}_{X/S})'$。
\end{SGAresult}
\begin{proof}证明是直接的。\end{proof}
\begin{SGAresult}{命题}{2}
设 $S$ 是\Term{characteristic}{特征}为 $p$ 的\Term{prescheme}{预概形}，$X$ 是 $S$-\Term{prescheme}{预概形}。

a) \Term{relative-frobenius}{相对 Frobenius 态射} $\mathrm{Fr}_{X/S}:X\to X^{(p/S)}$ 是\Term{integral-morphism}{整}\Term{surjective}{满射}且\Term{radicial}{根性}的（因而是\Term{universal-homeomorphism}{万有同胚}）。
\SourcePage{458}
b) 若 $X\to S$ 是\Term{open-immersion}{开浸入}，则 $X^{(p)}=X$ 且 $\mathrm{Fr}_{X/S}=\operatorname{id}_X$。若 $X\to S$ 是\Term{closed-immersion}{闭浸入}，将 $X$ 认同于 $S$ 的\Term{closed-ringed-subspace}{闭赋环子空间} $(|X|,\mathcal O_S/\mathcal J)$，则
\[
X^{(p)}=(|X|,\mathcal O_S/\mathcal J^p),\qquad
\mathrm{Fr}_{X/S}=(\operatorname{id}_{|X|},\varphi),
\]
其中 $\varphi:\mathcal O_S/\mathcal J^p\to\mathcal O_S/\mathcal J$ 是显然的\Term{augmentation}{增广}。

c) 假设 $X$ 在 $S$ 上\Term{locally-finite-presentation}{局部有限表示}。

1) 下列三个条件\Term{equivalence}{等价}：

(i) $X$ 在 $S$ 上\Term{unramified}{非分歧}；

(ii) $\mathrm{Fr}_{X/S}$ \Term{unramified}{非分歧}；

(iii) $\mathrm{Fr}_{X/S}$ 是\Term{monomorphism}{单态射}。

2) 同样，下列三个条件\Term{equivalence}{等价}：

(i') $X$ 在 $S$ 上\Term{etale}{平展}；

(ii') $\mathrm{Fr}_{X/S}$ \Term{etale}{平展}；

(iii') $\mathrm{Fr}_{X/S}$ 是\Term{isomorphism}{同构}。
\end{SGAresult}
\begin{proof}
先证 a)。复合 $\pi_{X/S}\circ\mathrm{Fr}_{X/S}=\mathrm{fr}_X$ 显然是\Term{integral-morphism}{整}\Term{surjective}{满射}且\Term{radicial}{根性}的，所以 $\mathrm{Fr}_{X/S}$ 是\Term{radicial}{根性}的（EGA I 3.5.6 (ii)）。此外，$\pi_{X/S}$ 由 $\mathrm{fr}_S$ 沿 $X\to S$ \Term{base-change}{基变换}得到，故它是\Term{separated}{分离}且\Term{radicial}{根性}的。因此 $\mathrm{Fr}_{X/S}$ 是\Term{integral-morphism}{整}的（EGA II 6.1.5 (v)）并且\Term{surjective}{满射}。

b) 的证明显然，只需注意：对 $\mathcal O_S$ 的任意\Term{ideal}{理想} $\mathcal J$，有 $\mathcal J\mathcal O_{(S,\mathrm{fr}_S)}=\mathcal J^p$。

证明 c)：我们知道，一个\Term{morphism}{态射}是\Term{monomorphism}{单态射}（相应地是\Term{isomorphism}{同构}）的充要条件是它\Term{radicial}{根性}且\Term{unramified}{非分歧}（相应地是\Term{radicial}{根性}、\Term{etale}{平展}且\Term{surjective}{满射}）。
\SourcePage{459}
故 (ii) 与 (iii) \Term{equivalence}{等价}，(ii') 与 (iii') \Term{equivalence}{等价}。还须证明 (i) 与 (ii)、(i') 与 (ii') 分别\Term{equivalence}{等价}。

为此，使用\Term{unramified}{非分歧}的\Term{differential}{微分}判别法（S.G.A. I 3），以及\Term{canonical}{典范}\Term{exact-sequence}{正合列}
\[
\mathrm{Fr}_{X/S}^{*}(\Omega^1_{X^{(p)}/S})\longrightarrow
\Omega^1_{X/S}\longrightarrow\Omega^1_{X/X^{(p)}}\longrightarrow0
\]
（S.G.A. II 4）给出的\Term{isomorphism}{同构} $\Omega^1_{X/S}\simeq\Omega^1_{X/X^{(p)}}$。确实，上列的第一个\Term{arrow}{箭头}显然为零，因为 $\Omega^1_{X^{(p)}/S}$ 由 $\mathcal O_{X^{(p)}}$ 的\Term{section}{截面}的\Term{differential}{微分}生成，而 $d_{X/S}$ 在 $(\mathcal O_X)^p$ 以及 $\mathcal O_S$ 上均为零。于是 (i) 与 (ii) \Term{equivalence}{等价}。

现设 $X$ 在 $S$ 上\Term{etale}{平展}，则 $X^{(p)}$ 亦然（\Term{base-change}{基变换}），因而 $\mathrm{Fr}_{X/S}:X\to X^{(p)}$ 是\Term{etale}{平展}的。反之，设 $\mathrm{Fr}_{X/S}$ \Term{etale}{平展}；于是它\Term{unramified}{非分歧}，故 $X$ 在 $S$ 上\Term{unramified}{非分歧}。换言之，结构\Term{morphism}{态射} $g:X\to S$ 在 $X$ 上\Term{local-scheme}{局部}地（对于\Term{etale-topology}{平展拓扑}）是\Term{closed-immersion}{闭浸入}（S.G.A. I 推论 7.8）。因此可假设 $g$ 是\Term{finite-presentation}{有限表示}的\Term{closed-immersion}{闭浸入}，只需说明：若 $\mathrm{Fr}_{X/S}$ 是\Term{isomorphism}{同构}，则 $g$ 是\Term{open-immersion}{开浸入}。此时 $X$ 由\Term{finite-type}{有限型}\Term{ideal}{理想} $\mathcal J$ 定义，即 $X=(|X|,\mathcal O_S/\mathcal J)$；由 b)，
\[
X^{(p)}=(|X|,\mathcal O_S/\mathcal J^p),\qquad
\mathrm{Fr}_{X/S}=(\operatorname{id}_{|X|},\mathcal O_S/\mathcal J^p\longrightarrow\mathcal O_S/\mathcal J).
\]
$\mathrm{Fr}_{X/S}$ 是\Term{isomorphism}{同构}\Term{equivalence}{等价}于 $\mathcal J=\mathcal J^p$。由于 $\mathcal J$ 为\Term{finite-type}{有限型}，\Term{nakayama-lemma}{Nakayama 引理}给出：对每个 $s\in S$，$\mathcal J_s=0$ 或 $\mathcal J_s=\mathcal O_{S,s}$，结论由此得证。
\end{proof}
\SourcePage{460}
\textbf{3. 对模的\Term{action}{作用}；例子。}

设 $S$ 是\Term{characteristic}{特征}为 $p$ 的\Term{prescheme}{预概形}，$X$ 是 $S$-\Term{prescheme}{预概形}。对任意 $\mathcal O_X$-\Term{module}{模} $\mathcal E$（Zariski \Term{topology}{拓扑}的\Term{sheaf}{层}），置
\[
\mathcal E^{(p/S)}=\pi_{X/S}^{*}(\mathcal E)
=\mathcal E\otimes_{\mathcal O_X}\mathcal O_{X^{(p)}}
=\mathcal E\otimes_{\mathcal O_S}\mathcal O_{(S,\mathrm{fr}_S)}.
\]
这是一个 $\mathcal O_{X^{(p)}}$-\Term{module}{模}，且有
\[
\mathcal E^{(p/X)}\simeq\mathrm{Fr}_{X/S}^{*}(\mathcal E^{(p/S)}).
\]
最有意义的情形是 $X=S$；此时将 $\mathcal E^{(p/S)}$ 写成 $\mathcal E^{(p)}$ 不会混淆，且 $\mathcal E^{(p)}=\mathrm{fr}_S^{*}(\mathcal E)$。\Term{adjunction}{伴随}\Term{morphism}{态射} $\mathcal E\to\mathrm{fr}_{S*}\mathrm{fr}_S^{*}(\mathcal E)$ 是阿贝尔\Term{sheaf}{层}的\Term{morphism}{态射} $\mathcal E\to\mathcal E^{(p)}$；它不是 $\mathcal O_S$-\Term{linear}{线性}的，而是“$p$-\Term{linear}{线性}的”（其含义显然），并且显然是从 $\mathcal E$ 到 $\mathcal O_S$-\Term{module}{模}的泛 $p$-\Term{linear}{线性}\Term{morphism}{态射}。

显然，$\mathcal E^{(p)}$ 对 $\mathcal E$ 具有\Term{functoriality}{函子性}，\Term{functor}{函子} $\mathcal E\mapsto\mathcal E^{(p)}$ 是\Term{right-exact}{右正合}的（一般并不\Term{exact}{正合}）；它与\Term{tensor-product}{张量积}交换：
\[
(\mathcal E\otimes_{\mathcal O_X}\mathcal F)^{(p)}
\simeq\mathcal E^{(p)}\otimes_{\mathcal O_{X^{(p)}}}\mathcal F^{(p)},
\]
并保持\Term{quasi-coherent}{拟凝聚}、\Term{finite-type}{有限型}、\Term{finite-presentation}{有限表示}、\Term{flat}{平坦}、\Term{locally-free}{局部自由}等性质。特别地，若 $\mathcal L$ 是 $S$ 上的\Term{invertible-sheaf}{可逆层}，则立即可见 $\mathcal L^{(p)}$ \Term{canonical}{典范}\Term{isomorphism}{同构}于 $\mathcal L^{\otimes p}$，因为由\Term{section}{截面}的 $p$ 次\Term{tensor-power}{张量幂}定义的\Term{morphism}{态射} $\mathcal L\to\mathcal L^{\otimes p}$ 是泛 $p$-\Term{linear}{线性}的。
\SourcePage{461}
若 $\mathcal A$ 是 $\mathcal O_S$-\Term{algebra}{代数}，则 $\mathcal A^{(p)}$ 带有由 $\mathcal A$ \Term{derive}{导出}的 $\mathcal O_S$-\Term{algebraic-structure}{代数结构}。假设 $\mathcal A$ \Term{quasi-coherent}{拟凝聚}（相应地\Term{graded}{分次}且\Term{quasi-coherent}{拟凝聚}），则 $\mathcal A^{(p)}$ 亦然，且
\[
\operatorname{Spec}(\mathcal A^{(p)})
=\operatorname{Spec}(\mathcal A\otimes_{\mathcal O_S}\mathcal O_{(S,\mathrm{fr}_S)})
\simeq\operatorname{Spec}(\mathcal A)\times_S(S,\mathrm{fr}_S)
=\operatorname{Spec}(\mathcal A)^{(p)},
\]
相应地，
\[
\operatorname{Proj}(\mathcal A^{(p)})
=\operatorname{Proj}(\mathcal A\otimes_{\mathcal O_S}\mathcal O_{(S,\mathrm{fr}_S)})
\simeq\operatorname{Proj}(\mathcal A)\times_S(S,\mathrm{fr}_S)
=\operatorname{Proj}(\mathcal A)^{(p)}.
\]
\Term{projection}{投影} $\pi_{\operatorname{Spec}(\mathcal A)/S}:\operatorname{Spec}(\mathcal A)^{(p)}\to\operatorname{Spec}(\mathcal A)$（相应地 $\pi_{\operatorname{Proj}(\mathcal A)/S}$）就是由 $\mathcal A\to\mathcal A^{(p)}$ \Term{derive}{导出}的\Term{morphism}{态射}。另一方面，$\mathrm{Fr}_{\operatorname{Spec}(\mathcal A)/S}$（相应地 $\mathrm{Fr}_{\operatorname{Proj}(\mathcal A)/S}$）来自 $\mathcal O_S$-\Term{algebra}{代数}\Term{homomorphism}{同态}
\[
\mathcal A^{(p)}=\mathcal A\otimes_{\mathcal O_S}\mathcal O_{(S,\mathrm{fr}_S)}\longrightarrow\mathcal A,
\qquad a\otimes f\longmapsto a^pf,
\]
这里 $a$、$f$ 分别为 $\mathcal A$、$\mathcal O_S$ 的\Term{section}{截面}。只需注意：此\Term{homomorphism}{同态}与右方的 $\mathcal A\to\mathcal A^{(p)}$ 复合，就是 $\mathcal A$ 中的 $p$ 次幂运算。

考虑 $\mathcal A$ 为\Term{quasi-coherent}{拟凝聚} $\mathcal O_S$-\Term{module}{模} $\mathcal E$ 的\Term{symmetric-algebra}{对称代数} $\operatorname{Sym}(\mathcal E)$ 的特例。显然 $\mathcal A^{(p)}\simeq\operatorname{Sym}(\mathcal E^{(p)})$，所以
\[
\begin{aligned}
\mathbf V(\mathcal E^{(p)})&\simeq\operatorname{Spec}(\mathcal A^{(p)})\simeq\operatorname{Spec}(\mathcal A)^{(p)}=\mathbf V(\mathcal E)^{(p)},\\
\mathbf P(\mathcal E^{(p)})&\simeq\operatorname{Proj}(\mathcal A^{(p)})\simeq\operatorname{Proj}(\mathcal A)^{(p)}=\mathbf P(\mathcal E)^{(p)}.
\end{aligned}
\]
至于 $\mathrm{Fr}_{\mathbf V(\mathcal E)/S}$（相应地 $\mathrm{Fr}_{\mathbf P(\mathcal E)/S}$），则\Term{resolution}{分解}为
\[
\mathbf V(\mathcal E)\longrightarrow\mathbf V(\operatorname{Sym}^p(\mathcal E))\longrightarrow\mathbf V(\mathcal E^{(p)})
\]
（相应地 $\mathbf P(\mathcal E)\to\mathbf P(\operatorname{Sym}^p(\mathcal E))\to\mathbf P(\mathcal E^{(p)})$）。
\SourcePage{462}
其中第一个\Term{arrow}{箭头}由来自 $\operatorname{Sym}^p(\mathcal E)\to\operatorname{Sym}(\mathcal E)$ 的\Term{canonical}{典范}\Term{morphism}{态射}
\[
\operatorname{Sym}(\operatorname{Sym}^p(\mathcal E))\longrightarrow\operatorname{Sym}(\mathcal E)
\]
定义；第二个\Term{arrow}{箭头}由 $\mathcal O_S$-\Term{module}{模}的\Term{morphism}{态射} $\mathcal E^{(p)}\to\operatorname{Sym}^p(\mathcal E)$ 定义，它来自取\Term{section}{截面}的 $p$ 次\Term{symmetric-power}{对称幂}所给出的 $p$-\Term{linear}{线性}\Term{morphism}{态射} $\mathcal E\to\operatorname{Sym}^p(\mathcal E)$。当 $X=\mathbf P(\mathcal E)$ 时，计算基本\Term{sheaf}{层} $\mathcal O_{X^{(p)}}(1)$ 得
\[
\mathcal O_{X^{(p)}}(1)
=(\operatorname{Sym}(\mathcal E)^{(p)}(1))^{\sim}
=(\operatorname{Sym}(\mathcal E)(1)^{(p)})^{\sim}
\simeq\mathcal O_X(1)^{(p/S)}.
\]
当 $\mathcal E=\mathcal O_S^n$ 时，$\operatorname{Sym}(\mathcal E)=\mathcal O_S[T_1,\ldots,T_n]$，且 $X=\mathbf V(\mathcal E)=\mathbf S[T_1,\ldots,T_n]$（相应地 $X=\mathbf P(\mathcal E)=\mathbf P_S^{n-1}$）。显然 $\mathcal E^{(p)}\simeq\mathcal O_S^n=\mathcal E$，故 $X^{(p)}=X$。\Term{relative-frobenius}{Frobenius 态射} $\mathrm{Fr}_{X/S}$ 来自\Term{algebra}{代数} $\mathcal O_S[T_1,\ldots,T_n]$ 的\Term{endomorphism}{自同态} $T_i\mapsto T_i^p$（$i=1,\ldots,n$）。

现考虑 $S$ 上\Term{finite-type}{有限型}\Term{affine}{仿射}（相应地\Term{projective}{射影}）\Term{prescheme}{预概形} $X$，由\Term{ideal}{理想}（相应地\Term{homogeneous-ideal}{齐次理想}）$\mathcal J\subset\mathcal O_S[T_1,\ldots,T_n]$ 定义。于是 $X=\operatorname{Spec}(\mathcal A)$（相应地 $X=\operatorname{Proj}(\mathcal A)$），其中 $\mathcal A=\mathcal O_S[T_1,\ldots,T_n]/\mathcal J$。可见
\[
\mathcal A^{(p)}\simeq\mathcal O_S[T_1,\ldots,T_n]/\mathcal J',
\]
这里 $\mathcal J'$ 是 $\mathcal J^{(p)}=\mathcal J\otimes_{\mathcal O_S}\mathcal O_{(S,\mathrm{fr}_S)}$ 在 $\mathcal O_S[T_1,\ldots,T_n]$ 中的像；$\mathcal J'$ 由 $\mathcal J$ 中的\Term{polynomial}{多项式}的系数取 $p$ 次幂得到。于是 $X^{(p)}=\operatorname{Spec}(\mathcal A^{(p)})$（相应地 $\operatorname{Proj}(\mathcal A^{(p)})$），而\Term{relative-frobenius}{Frobenius 态射} $\mathrm{Fr}_{X/S}:X\to X^{(p)}$ 由“坐标” $T_i$ 的 $p$ 次幂运算定义（该运算确定\Term{morphism}{态射} $\mathcal A^{(p)}\to\mathcal A$）。
\SourcePage{463}
\textbf{4. 迭代 \Term{frobenius-morphism}{Frobenius 态射}；\Term{perfect-prescheme}{完美预概形}。}

对每个整数 $\nu\geq0$，迭代\Term{endomorphism}{自同态} $\mathrm{fr}_X^\nu$ 对\Term{characteristic}{特征}为 $p$ 的\Term{prescheme}{预概形} $X$ 具有\Term{functoriality}{函子性}。若 $S$ 是\Term{characteristic}{特征}为 $p$ 的\Term{prescheme}{预概形}，$X$ 是 $S$-\Term{prescheme}{预概形}，则 $\mathrm{fr}_X^\nu$ 唯一地\Term{resolution}{分解}为 $S$-\Term{morphism}{态射}
\[
\mathrm{Fr}_{X/S}^{\nu}:X\longrightarrow X^{(p^\nu/S)},
\qquad X^{(p^\nu/S)}=X\times_S(S,\mathrm{fr}_S^\nu),
\]
再接\Term{projection}{投影} $X^{(p^\nu/S)}\to X$（参见第 2 小节）。\Term{morphism}{态射} $\mathrm{Fr}_{X/S}^{\nu}$ 对 $S$-\Term{prescheme}{预概形} $X$ 具有\Term{functoriality}{函子性}，并与底 $S$ 的\Term{base-change}{基变换}相容（参见第 2 小节命题 1 b)）。此外，显然有
\[
X^{(p^\nu/S)}=X^{(p^\nu)}\simeq(X^{(p^{\nu-1})})^{(p)},
\qquad
\mathrm{Fr}_{X/S}^{\nu}=(\mathrm{Fr}_{X/S}^{\nu-1})^{(p)}\circ\mathrm{Fr}_{X/S},
\]
其中使用上述\Term{canonical}{典范}\Term{isomorphism}{同构}进行认同。
\begin{SGAresult}{定义}{4}
若\Term{characteristic}{特征}为 $p$ 的\Term{prescheme}{预概形} $S$ 的 $\mathrm{fr}_S$ 是 $S$ 的\Term{automorphism}{自同构}，则称 $S$ 为\Term{perfect-prescheme}{完美预概形}。
\end{SGAresult}
若 $S$ 是\Term{characteristic}{特征}为 $p$ 的\Term{perfect-prescheme}{完美预概形}，$X$ 是 $S$-\Term{prescheme}{预概形}，则立即可见，$X^{(p/S)}$ 可认同于以 $\mathrm{fr}_S^{-1}\circ g:X\to S$ 为结构\Term{morphism}{态射}的 $X$，其中 $g$ 为 $X$ 到 $S$ 的结构\Term{morphism}{态射}；在此认同下 $\mathrm{Fr}_{X/S}=\mathrm{fr}_X$。

以后会遇到的\Term{perfect-prescheme}{完美预概形}的一个例子是 $S=\operatorname{Spec}(\mathbb F_q)$，其中 $q=p^\nu$ 是\Term{characteristic}{特征} $p$ 的幂；此时 $\mathrm{fr}_S^\nu=\operatorname{id}_S$。若 $X$ 是 $S$-\Term{prescheme}{预概形}，可见
\SourcePage{464}
$X^{(p^\nu/S)}=X$ 且 $\mathrm{Fr}_{X/S}^{\nu}=\mathrm{fr}_X^\nu$。特别地，$\nu=1$ 时得到 $S=\operatorname{Spec}(\mathbb F_p)$，其 $\mathrm{fr}_S=\operatorname{id}_S$；$S$-\Term{prescheme}{预概形}就是\Term{characteristic}{特征}为 $p$ 的\Term{prescheme}{预概形}，在这一底 $S$ 上，$X^{(p/S)}=X$ 且 $\mathrm{Fr}_{X/S}=\mathrm{fr}_X$。

\SGAsection{2}{Frobenius 对应}
\textbf{1. 定义。}

设 $X$ 是\Term{characteristic}{特征}为 $p$ 的\Term{prescheme}{预概形}，$F$ 是 $X_{\text{ét}}$ 上的\Term{set}{集合}\Term{sheaf}{层}。对每个\Term{etale}{平展} $X$-\Term{prescheme}{预概形} $U$，有
\[
(\mathrm{fr}_X)_*(F)(U)=F(\mathrm{fr}_X^{-1}(U))=F(U^{(p/X)}),
\]
而\Term{relative-frobenius}{Frobenius 态射} $\mathrm{Fr}_{U/X}:U\to U^{(p/X)}$ 定义映射
\[
F(\mathrm{Fr}_{U/X}):\mathrm{fr}_{X*}(F)(U)\longrightarrow F(U),
\]
它显然对 $U$ 具有\Term{functoriality}{函子性}（因为 $\mathrm{Fr}_{U/X}$ 对 $U$ 具有\Term{functoriality}{函子性}）。由此得到 $X_{\text{ét}}$ 上的\Term{sheaf}{层}\Term{morphism}{态射}
\[
(\mathrm{fr}_X)_*(F)\longrightarrow F.
\]
当 $U$ 在 $X$ 上\Term{etale}{平展}时，$\mathrm{Fr}_{U/X}$ 是\Term{isomorphism}{同构}（第 1 节第 2 小节命题 2 c)），故 $F(\mathrm{Fr}_{U/X})$ 亦然。因此上述\Term{morphism}{态射} $(\mathrm{fr}_X)_*(F)\to F$ 是\Term{isomorphism}{同构}；若 $F$ 带有\Term{algebraic-structure}{代数结构}（例如它是 $\Lambda$-\Term{module}{模}\Term{sheaf}{层}，$\Lambda$ 为给定的系数环），则此\Term{isomorphism}{同构}与该\Term{algebraic-structure}{代数结构}相容。此\Term{isomorphism}{同构}的逆
\[
F(\mathrm{Fr}_{/X})^{-1}:F\longrightarrow\mathrm{fr}_{X*}(F)
\]
借助 $(\mathrm{fr}_X)^*$ 与 $(\mathrm{fr}_X)_*$ 的\Term{left-adjoint}{左伴随}关系，定义了一个\Term{morphism}{态射}
\SourcePage{465}
\[
\mathrm{Fr}_{F/X}^{*}:(\mathrm{fr}_X)^*(F)\longrightarrow F.
\]
还可以进一步说：由于 $\mathrm{fr}_X$ 是\Term{universal-homeomorphism}{万有同胚}（第 1 节第 2 小节命题 2 a)），\Term{functor}{函子} $(\mathrm{fr}_X)^*$ 与 $(\mathrm{fr}_X)_*$ 互为\Term{quasi-inverse}{拟逆}，所以 $\mathrm{Fr}_{F/X}^{*}$ 与产生它的\Term{morphism}{态射} $F(\mathrm{Fr}_{/X})^{-1}$ 一样，也是\Term{isomorphism}{同构}。
\begin{SGAresult}{定义}{1}
给定\Term{characteristic}{特征}为 $p$ 的\Term{prescheme}{预概形} $X$ 以及 $X_{\text{ét}}$ 上的\Term{set}{集合}\Term{sheaf}{层} $F$，称由 $X$ 的\Term{frobenius-endomorphism}{Frobenius 自同态} $\mathrm{fr}_X$ 与\Term{isomorphism}{同构}
\[
\mathrm{Fr}_{F/X}^{*}=(F(\mathrm{Fr}_{/X})^{-1})^{\sharp}:
(\mathrm{fr}_X)^*(F)\longrightarrow F
\]
组成的二元组 $(\mathrm{fr}_X,\mathrm{Fr}_{F/X}^{*})$ 为 $(X,F)$ 上的 \Term{frobenius-correspondence}{Frobenius 对应}。
\end{SGAresult}
当然，对任意整数 $\nu\geq0$，同样可定义\Term{iterated-frobenius-correspondence}{迭代 Frobenius 对应} $(\mathrm{fr}_X^\nu,(\mathrm{Fr}_{F/X}^{*})^\nu)$；当 $\nu\geq1$ 时，\Term{isomorphism}{同构} $(\mathrm{Fr}_{F/X}^{*})^\nu:(\mathrm{fr}_X^\nu)^*(F)\to F$ 等于
\[
(\mathrm{Fr}_{F/X}^{*})^{\nu-1}\circ(\mathrm{fr}_X^{\nu-1})^*(\mathrm{Fr}_{F/X}^{*})
\]
也等于 $\mathrm{Fr}_{F/X}^{*}\circ\mathrm{fr}_X^*((\mathrm{Fr}_{F/X}^{*})^{\nu-1})$。

考察 $F$ 可由 $X$ 上\Term{etale}{平展}的\Term{prescheme}{预概形} $V$ 表示的特例。此时 $\mathrm{fr}_X^*(F)=V^{(p/X)}$，
\[
F(V)=\operatorname{Hom}_X(V,V),\quad
\mathrm{fr}_{X*}(F)(V)=F(V^{(p/X)})=\operatorname{Hom}_X(V^{(p/X)},V),
\]
而 $F(\mathrm{Fr}_{V/X})$ 是从 $\operatorname{Hom}_X(V^{(p/X)},V)$ 到 $\operatorname{Hom}_X(V,V)$ 的映射 $\varphi\mapsto\varphi\circ\mathrm{Fr}_{V/X}$。所以 $\mathrm{Fr}_{F/X}^{*}$ 可认同于 $V^{(p/X)}\to V$ 的\Term{morphism}{态射}，它是 $\operatorname{id}_V$ 在上述映射的逆映射 $F(\mathrm{Fr}_{V/X})^{-1}$ 下的像。这表明：对 $X$ 上\Term{etale}{平展}的 $F=V$，有
\[
\mathrm{Fr}_{V/X}^{*}=(\mathrm{Fr}_{V/X})^{-1}.
\]
\SourcePage{466}
注意，若 $F$ 是\Term{constant-sheaf}{常值层}，则显然 $\mathrm{fr}_X^*(F)=F$ 且 $\mathrm{Fr}_{F/X}^{*}=\operatorname{id}_F$。
\begin{SGAresult}{命题}{1}
a) \Term{isomorphism}{同构} $\mathrm{Fr}_{F/X}^{*}$ 对\Term{sheaf}{层} $F$ 具有\Term{functoriality}{函子性}；换言之，有自然\Term{isomorphism}{同构}
\[
\mathrm{Fr}_{/X}^{*}:\mathrm{fr}_X^*\longrightarrow\operatorname{id}_{\widetilde X_{\text{ét}}},
\]
其中 $\widetilde X_{\text{ét}}$ 为 $X_{\text{ét}}$ 上的\Term{set}{集合}\Term{sheaf}{层}的\Term{category}{范畴}。

b) $(X,F)$ 上的\Term{frobenius-correspondence}{Frobenius 对应} $(\mathrm{fr}_X,\mathrm{Fr}_{F/X}^{*})$ 与底 $X$ 的\Term{base-change}{基变换}相容。具体地，若 $\mathscr F$ 表示在 $(\mathrm{Sch})_p$ 上、由变动的\Term{characteristic}{特征}为 $p$ 的\Term{prescheme}{预概形}上的\Term{etale-topology}{平展拓扑}\Term{sheaf}{层}组成的\Term{fibered-category}{纤维范畴}，则\Term{functor}{函子} $\mathrm{fr}_X^*$（$X\in\operatorname{Ob}(\mathrm{Sch})_p$）定义一个 $(\mathrm{Sch})_p$-\Term{cartesian-functor}{笛卡尔函子} $\mathrm{fr}^*:\mathscr F\to\mathscr F$，而 $\mathrm{Fr}_{F/X}^{*}$（$X\in\operatorname{Ob}(\mathrm{Sch})_p$，$F\in\operatorname{Ob}\widetilde X_{\text{ét}}$）定义一个 $(\mathrm{Sch})_p$-自然\Term{morphism}{态射}
\[
\mathrm{Fr}^{*}:\mathrm{fr}^{*}\longrightarrow\operatorname{id}_{\mathscr F},
\]
且它是\Term{isomorphism}{同构}。

c) 反之，这些性质刻画了 $\mathrm{Fr}^{*}$；即从 $\mathrm{fr}^{*}$ 到 $\operatorname{id}_{\mathscr F}$ 的 $(\mathrm{Sch})_p$-自然\Term{morphism}{态射}存在且唯一。
\end{SGAresult}
\begin{proof}
注意 b) 蕴含 a)，而 a) 完全显然。

证明 b)。设 $g:X'\to X$ 是\Term{characteristic}{特征}为 $p$ 的\Term{prescheme}{预概形}之间的\Term{morphism}{态射}，$F$ 是 $X_{\text{ét}}$ 上的\Term{sheaf}{层}。置 $F'=g^*(F)$，更一般地，用撇号表示 $X$ 上任意\Term{object}{对象}沿 $g$ 的\Term{inverse-image}{逆像}。由于 $g\circ\mathrm{fr}_{X'}=\mathrm{fr}_X\circ g$，有
\SourcePage{467}
\[
\mathrm{fr}_{X'}^{*}(F')=\mathrm{fr}_{X'}^{*}(g^*(F))
=g^*(\mathrm{fr}_X^*(F))=(\mathrm{fr}_X^*(F))'.
\]
故各 $\mathrm{fr}_X^*$ 确实定义 $(\mathrm{Sch})_p$ 上的\Term{cartesian-functor}{笛卡尔函子} $\mathrm{fr}^*$。还须证明 $\mathrm{Fr}_{F'/X'}^{*}=(\mathrm{Fr}_{F/X}^{*})'$。利用\Term{adjoint-functor}{伴随函子}的性质，可见这\Term{equivalence}{等价}于：$F'(\mathrm{Fr}_{/X'})^{-1}$ 经\Term{bijection}{双射}
\[
\begin{aligned}
\operatorname{Hom}_{X'}(F',\mathrm{fr}_{X'*}(F'))
&=\operatorname{Hom}_{X'}(g^*(F),\mathrm{fr}_{X'*}g^*(F))\\
&\xrightarrow{\sim}\operatorname{Hom}_X(F,g_*\mathrm{fr}_{X'*}g^*(F))
\end{aligned}
\]
（$g^*$ 与 $g_*$ \Term{adjunction}{伴随}）所得的像，等于 $F(\mathrm{Fr}_{/X})^{-1}$ 经由\Term{adjunction}{伴随}\Term{morphism}{态射} $\alpha:F\to g_*g^*(F)$ 定义的映射
\[
\begin{aligned}
\operatorname{Hom}_X(F,\mathrm{fr}_{X*}(F))
&\longrightarrow\operatorname{Hom}_X(F,\mathrm{fr}_{X*}g_*g^*(F))\\
&=\operatorname{Hom}_X(F,g_*\mathrm{fr}_{X'*}g^*(F))
\end{aligned}
\]
所得的像。换言之，须证明
\[
\alpha\circ F(\mathrm{Fr}_{/X})
=g_*(F'(\mathrm{Fr}_{/X'}))\circ\mathrm{fr}_{X*}(\alpha),
\]
即对每个在 $X$ 上\Term{etale}{平展}的\Term{prescheme}{预概形} $U$，下图交换：
\[
\begin{tikzcd}[column sep=large]
F(U^{(p/X)})\arrow[r,"F(\mathrm{Fr}_{U/X})"]\arrow[d]&F(U)\arrow[d]\\
F'(U'^{(p/X')})\arrow[r,"F'(\mathrm{Fr}_{U'/X'})"']&F'(U').
\end{tikzcd}
\]
其中竖\Term{arrow}{箭头}为\Term{canonical-map}{典范映射}；这直接来自 $\mathrm{Fr}_{U'/X'}=(\mathrm{Fr}_{U/X})'$（第 1 节第 2 小节命题 1 b)）。

为证明 c)，与 $(\mathrm{Fr}^{*})^{-1}$ 复合后，只需证明 $\operatorname{id}_{\mathscr F}$ 的唯一 $(\mathrm{Sch})_p$-\Term{endomorphism}{自同态}为恒等。设 $\varphi$ 为这样的\Term{endomorphism}{自同态}，$F$ 是\Term{characteristic}{特征}为 $p$ 的\Term{prescheme}{预概形} $X$ 上的\Term{sheaf}{层}，$U$ 是 $X$ 上\Term{etale}{平展}的\Term{prescheme}{预概形}。我们知道 $F(U)$ 可认同于 $\operatorname{Hom}_U(U,F|U)$，
\SourcePage{468}
而 $\varphi_F(U):F(U)\to F(U)$ 则认同于映射
\[
s\longmapsto\varphi_{F|U}\circ s=s\circ\varphi_U,
\qquad s\in\operatorname{Hom}_U(U,F|U).
\]
由于 $U_{\text{ét}}$ 的终\Term{sheaf}{层} $U$ 除恒等之外没有别的\Term{endomorphism}{自同态}，有 $\varphi_U=\operatorname{id}_U$，于是 $\varphi_F(U)$ 也是 $F(U)$ 的\Term{identity-map}{恒等映射}，从而 $\varphi=\operatorname{id}_{\mathscr F}$。
\end{proof}
当然，若 $F$ 带有\Term{algebraic-structure}{代数结构}，\Term{isomorphism}{同构} $\mathrm{Fr}_{F/X}^{*}$ 与该结构相容。例如，若 $F$ 是\Term{characteristic}{特征}为 $p$ 的\Term{prescheme}{预概形} $X$ 上的群\Term{sheaf}{层}，则 $\mathrm{Fr}_{F/X}^{*}$ 是群\Term{sheaf}{层}的\Term{isomorphism}{同构}。设 $\Lambda$ 为环，$F$ 是 $X$ 上的 $\Lambda$-\Term{module}{模}\Term{sheaf}{层}，则 $\mathrm{Fr}_{F/X}^{*}$ 是 $\Lambda$-\Term{module}{模}\Term{sheaf}{层}的\Term{isomorphism}{同构}。更一般地，对 $\Lambda$-\Term{module}{模}\Term{complex}{复形} $F^{\bullet}$，立即可定义\Term{isomorphism}{同构}
\[
\mathrm{Fr}_{F^{\bullet}/X}^{*}:\mathrm{fr}_X^*(F^{\bullet})\longrightarrow F^{\bullet},
\]
其在\Term{degree}{次数} $i$ 的\Term{set-component}{分量}等于 $\mathrm{Fr}_{F^i/X}^{*}$（由命题 1 a)）。此时二元组 $(\mathrm{fr}_X,\mathrm{Fr}_{F^{\bullet}/X}^{*})$ 确实是 $(X,F^{\bullet})$ 上通常意义的\Term{correspondence}{对应}；它对 $F^{\bullet}$ 具有\Term{functoriality}{函子性}，并与底 $X$ 的\Term{base-change}{基变换}相容（命题 1 a)、b)）。

\textbf{2. 对\Term{cohomology}{上同调}的\Term{action}{作用}。}
\begin{SGAresult}{命题}{2}
设 $X$ 是\Term{characteristic}{特征}为 $p$ 的\Term{prescheme}{预概形}。

a) 对 $X_{\text{ét}}$ 上任意\Term{set}{集合}\Term{sheaf}{层} $F$，由 $(\mathrm{fr}_X,\mathrm{Fr}_{F/X}^{*})$ 定义的 $H^0(X,F)$ 的\Term{endomorphism}{自同态}为恒等。它即复合
\[
H^0(X,F)\longrightarrow H^0(X,\mathrm{fr}_X^*(F))\longrightarrow H^0(X,F),
\]
其中第一个\Term{arrow}{箭头}是\Term{canonical-map}{典范映射}（来自 $H^0(X,F)$ 对 $X$ 的\Term{contravariant}{反变}性），第二个是 $H^0(X,\mathrm{Fr}_{F/X}^{*})$。
\SourcePage{469}
b) 对 $X_{\text{ét}}$ 上任意群\Term{sheaf}{层} $F$，由 $(\mathrm{fr}_X,\mathrm{Fr}_{F/X}^{*})$ 定义的 $H^1(X,F)$ 的\Term{endomorphism}{自同态}，即
\[
H^1(X,F)\longrightarrow H^1(X,\mathrm{fr}_X^*(F))\longrightarrow H^1(X,F),
\]
为恒等。

c) 对 $X_{\text{ét}}$ 上任意\Term{bounded-below}{左有界}的 $\Lambda$-\Term{module}{模}\Term{complex}{复形} $F$，由 $(\mathrm{fr}_X,\mathrm{Fr}_{F/X}^{*})$ 定义的 $\mathrm R\Gamma_X(F)$ 的\Term{endomorphism}{自同态}，即复合
\[
\mathrm R\Gamma_X(F)\longrightarrow\mathrm R\Gamma_X(\mathrm{fr}_X^*(F))\longrightarrow\mathrm R\Gamma_X(F),
\]
为恒等。
\end{SGAresult}
\begin{proof}
先证 a)，再由此\Term{pushout}{推出} b)、c)。以 $\varphi_F$ 表示所考察的 $H^0(X,F)$ 的\Term{endomorphism}{自同态}。显然它对 $F$ 具有\Term{functoriality}{函子性}，且 $\varphi_X=\operatorname{id}_{H^0(X,X)}$。而 $H^0(X,F)$ 的每个元素 $s$ 可认同于\Term{sheaf}{层}\Term{morphism}{态射} $s:X\to F$，使得 $s$ 为 $H^0(X,X)$ 的唯一元素在 $H^0(X,s)$ 下的像。根据上述说明，
\[
\varphi_F\circ H^0(X,s)=H^0(X,s)\circ\varphi_X=H^0(X,s),
\]
故对所有 $s$，$\varphi_F(s)=s$，于是 $\varphi_F=\operatorname{id}_{H^0(X,F)}$。

注意，$\varphi_F$ 也等于复合
\[
H^0(X,F)\longrightarrow H^0(X,\mathrm{fr}_{X*}(F))\longrightarrow H^0(X,F),
\]
其中第一个\Term{arrow}{箭头}是 $H^0(X,F(\mathrm{Fr}_{/X})^{-1})$，第二个是\Term{canonical-map}{典范映射}（也就是恒等）。同样，b)、c) 中考察的\Term{endomorphism}{自同态}分别\Term{resolution}{分解}为
\[
H^1(X,F)\longrightarrow H^1(X,\mathrm{fr}_{X*}(F))\longrightarrow H^1(X,F),
\]
以及
\[
\mathrm R\Gamma_X(F)\longrightarrow\mathrm R\Gamma_X(\mathrm{fr}_{X*}(F))\longrightarrow\mathrm R\Gamma_X(F).
\]
为证明 b)，取一个将 $F$ \Term{embedding}{嵌入}群\Term{sheaf}{层} $G$ 的\Term{monomorphism}{单态射} $F\to G$，使 $H^1(X,G)=0$。则 $\mathrm{fr}_{X*}(F)\to\mathrm{fr}_{X*}(G)$ 为\Term{monomorphism}{单态射}，且 $H^1(X,\mathrm{fr}_{X*}(G))=0$。引入商\Term{homogeneous-space}{齐性空间}的\Term{sheaf}{层} $C=G/F$。
\SourcePage{470}
由于 $\mathrm{fr}_X$ 是\Term{integral-morphism}{整}\Term{surjective}{满射}且\Term{radicial}{根性}的，\Term{functor}{函子} $\mathrm{fr}_X^*$ 与 $\mathrm{fr}_{X*}$ 是互为\Term{quasi-inverse}{拟逆}的\Term{equivalence}{范畴等价}（SGAA VIII \Term{theorem}{定理} 1.1），因而
\[
\mathrm{fr}_{X*}(C)\simeq\mathrm{fr}_{X*}(G)/\mathrm{fr}_{X*}(F).
\]
\Term{cohomology}{上同调}\Term{exact-sequence}{正合列}（SGAA XII 命题 3.1）表明，所考察的 $H^1(X,F)$ 的\Term{endomorphism}{自同态}是由 $\varphi_C:H^0(X,C)\to H^0(X,C)$ 取商得到的；后者由 a) 为恒等，故该\Term{endomorphism}{自同态}为 $\operatorname{id}_{H^1(X,F)}$。

c) 的证明类似。取一个各次均\Term{injective-module}{内射}、在 $\Lambda$-\Term{module}{模}上构造的\Term{derived-category}{导出范畴} $D^+$ 中与 $F$ \Term{isomorphism}{同构}的\Term{complex}{复形} $I$。由于 $\mathrm{fr}_{X*}$ 是\Term{equivalence}{范畴等价}，$\mathrm{fr}_{X*}(I)$ 各次均\Term{injective-module}{内射}，且在 $D^+$ 中与 $\mathrm{fr}_{X*}(F)$ \Term{isomorphism}{同构}。于是\Term{separable}{可分}别将 $\mathrm R\Gamma_X(F)$、$\mathrm R\Gamma_X(\mathrm{fr}_{X*}(F))$ 认同于 $H^0(X,I)$、$H^0(X,\mathrm{fr}_{X*}(I))$。所考察的 $\mathrm R\Gamma_X(F)$ 的\Term{endomorphism}{自同态}便认同于 $\varphi_I$；它为恒等，所以该\Term{endomorphism}{自同态}是 $\operatorname{id}_{\mathrm R\Gamma_X(F)}$。
\end{proof}
\textbf{3. 乘积的行为。}

设 $X$、$X'$ 是\Term{characteristic}{特征}为 $p$ 的\Term{prescheme}{预概形}；其乘积 $X\times X'$ 就是\Term{fiber-product}{纤维积} $X\times_eX'$，其中 $e=\operatorname{Spec}(\mathbb F_p)$。显然
\[
\mathrm{fr}_{X\times X'}=\mathrm{fr}_X\times\mathrm{fr}_{X'}
=(\mathrm{fr}_X\times\operatorname{id}_{X'})\circ(\operatorname{id}_X\times\mathrm{fr}_{X'})
=(\operatorname{id}_X\times\mathrm{fr}_{X'})\circ(\mathrm{fr}_X\times\operatorname{id}_{X'}).
\]
立即\Term{pushout}{推出}：对任意 $X$-\Term{prescheme}{预概形} $Y$ 及 $X'$-\Term{prescheme}{预概形} $Y'$，有\Term{canonical}{典范}\Term{isomorphism}{同构}
\[
(Y\times Y')^{(p/X\times X')}\simeq Y^{(p/X)}\times Y'^{(p/X')};
\]
在此\Term{isomorphism}{同构}下可写为
\SourcePage{471}
$\mathrm{Fr}_{Y\times Y'/X\times X'}=\mathrm{Fr}_{Y/X}\times\mathrm{Fr}_{Y'/X'}$。
\begin{SGAresult}{命题}{3}
设 $X$、$X'$ 是\Term{characteristic}{特征}为 $p$ 的\Term{prescheme}{预概形}，$\Lambda$ 为\Term{commutative-ring}{交换环}，$F$ 是 $X_{\text{ét}}$ 上的 $\Lambda$-\Term{module}{模}\Term{sheaf}{层}，$F'$ 是 $X'_{\text{ét}}$ 上的 $\Lambda$-\Term{module}{模}\Term{sheaf}{层}。则
\[
\begin{aligned}
\mathrm{Fr}_{F\boxtimes_\Lambda F'/X\times X'}^{*}
&=\mathrm{Fr}_{F/X}^{*}\boxtimes\mathrm{Fr}_{F'/X'}^{*}\\
&=(\mathrm{Fr}_{F/X}^{*}\boxtimes\operatorname{id}_{\mathrm{fr}_{X'}^*(F')})
\circ(\operatorname{id}_F\boxtimes\mathrm{Fr}_{F'/X'}^{*})\\
&=(\operatorname{id}_{\mathrm{fr}_X^*(F)}\boxtimes\mathrm{Fr}_{F'/X'}^{*})
\circ(\mathrm{Fr}_{F/X}^{*}\boxtimes\operatorname{id}_{F'}).
\end{aligned}
\]
此处利用\Term{canonical}{典范}\Term{isomorphism}{同构}，将 $\mathrm{fr}_{X\times X'}^*(F\boxtimes_\Lambda F')$ 认同于\Term{tensor-product}{张量积} $\mathrm{fr}_X^*(F)\boxtimes_\Lambda\mathrm{fr}_{X'}^*(F')$。
\end{SGAresult}
\begin{proof}
第一个等式的两边均对 $F$、$F'$ 具有\Term{functoriality}{函子性}，且与底 $X$ 或 $X'$ 的\Term{base-change}{基变换}相容；它们也与 $F$ 或 $F'$ 的\Term{inductive-limit}{归纳极限}相容，因为\Term{tensor-product}{张量积}与\Term{inductive-limit}{归纳极限}交换，并且显然
\[
\mathrm{Fr}_{(\varinjlim F_i)/X}^{*}=\varinjlim\mathrm{Fr}_{F_i/X}^{*}.
\]
因而\Term{reducible}{可约}化到 $F=\Lambda_X$、$F'=\Lambda_{X'}$，此时结论显然。
\end{proof}
\begin{SGAresult}{推论}{}
在命题 3 的数据下，分别由
\[
(\operatorname{id}_X\times\mathrm{fr}_{X'},\operatorname{id}_F\boxtimes\mathrm{Fr}_{F'/X'}^{*}),\qquad
(\mathrm{fr}_X\times\operatorname{id}_{X'},\mathrm{Fr}_{F/X}^{*}\boxtimes\operatorname{id}_{F'})
\]
定义的 $\mathrm R\Gamma_{X\times X'}(F\boxtimes F')$ 的\Term{endomorphism}{自同态}互为逆\Term{isomorphism}{同构}。
\end{SGAresult}
\begin{proof}
将这些\Term{endomorphism}{自同态}分别记为 $\varphi'$、$\varphi$，例如证明 $\varphi\circ\varphi'=\operatorname{id}$。这由下图得出：
\SourcePage{472}
\[
\resizebox{\textwidth}{!}{%
\begin{tikzcd}[ampersand replacement=\&,column sep=small,row sep=large,cramped]
\mathrm R\Gamma_{X\times X'}(F\boxtimes F')
\arrow[r,"\operatorname{id}_X\times\mathrm{fr}_{X'}"]
\arrow[dr,"\mathrm{fr}_{X\times X'}"']\&
\mathrm R\Gamma_{X\times X'}(F\boxtimes\mathrm{fr}_{X'}^*(F'))
\arrow[r,"{\mathrm R\Gamma(\operatorname{id}_F\boxtimes\mathrm{Fr}_{F'/X'}^*)}"]
\arrow[d,"\mathrm{fr}_X\times\operatorname{id}_{X'}"]\&
\mathrm R\Gamma_{X\times X'}(F\boxtimes F')
\arrow[d,"\mathrm{fr}_X\times\operatorname{id}_{X'}"]\\
\&\mathrm R\Gamma_{X\times X'}(\mathrm{fr}_{X\times X'}^*(F\boxtimes F'))
\arrow[r,"{\mathrm R\Gamma(\operatorname{id}_{\mathrm{fr}_X^*(F)}\boxtimes\mathrm{Fr}_{F'/X'}^*)}"]
\arrow[dr,"{\mathrm R\Gamma(\mathrm{Fr}_{F\boxtimes F'/X\times X'}^*)}"']\&
\mathrm R\Gamma_{X\times X'}(\mathrm{fr}_X^*(F)\boxtimes F')
\arrow[d,"{\mathrm R\Gamma(\mathrm{Fr}_{F/X}^*\boxtimes\operatorname{id}_{F'})}"]\\
\&\&\mathrm R\Gamma_{X\times X'}(F\boxtimes F').
\end{tikzcd}%
}
\]
它交换：上方三角形的\Term{commutativity}{交换性}来自 $\mathrm{fr}_{X\times X'}=(\mathrm{fr}_X\times\operatorname{id}_{X'})\circ(\operatorname{id}_X\times\mathrm{fr}_{X'})$；下方三角形的\Term{commutativity}{交换性}来自
\[
\mathrm{Fr}_{F\boxtimes F'/X\times X'}^{*}
=(\mathrm{Fr}_{F/X}^{*}\boxtimes\operatorname{id}_{F'})
\circ(\operatorname{id}_{\mathrm{fr}_X^*(F)}\boxtimes\mathrm{Fr}_{F'/X'}^{*}).
\]
由命题 2 c)（第 2 小节），两个斜\Term{arrow}{箭头}的复合为恒等。类似地，交换 $(X,F)$ 与 $(X',F')$ 的地位，得到 $\varphi'\circ\varphi=\operatorname{id}$。
\end{proof}
我们将把它用于如下特例：$X'=\bar e=\operatorname{Spec}(\overline{\mathbb F}_p)$，其中 $\overline{\mathbb F}_p$ 是 $\mathbb F_p$ 的\Term{algebraic-closure}{代数闭包}，$F'=\Lambda_{\bar e}$ 是以 $\Lambda$ 为\Term{fiber}{纤维}的\Term{constant-sheaf}{常值层}。此时置 $\bar X=X\times\bar e=X\otimes_{\mathbb F_p}\overline{\mathbb F}_p$，$\bar F=F\boxtimes_\Lambda\Lambda_{\bar e}$，它是 $F$ 沿\Term{projection}{投影} $\bar X\to X$ 的\Term{inverse-image}{逆像}。\Term{frobenius-endomorphism}{Frobenius 自同态} $\mathrm{fr}_{\bar e}$ 可认同于\Term{galois-group}{Galois 群} $\pi_e=\operatorname{Gal}(\overline{\mathbb F}_p/\mathbb F_p)$ 的\Term{canonical}{典范}\Term{topological-generator}{生成元} $f:\lambda\mapsto\lambda^p$（$\lambda\in\overline{\mathbb F}_p$；$f$ 确定\Term{isomorphism}{同构} $\widehat{\mathbb Z}\xrightarrow{\sim}\pi_e$），且 $\mathrm{Fr}_{\Lambda_{\bar e}/\bar e}^{*}=\operatorname{id}_{\Lambda_{\bar e}}$，因为 $\Lambda_{\bar e}$ 是\Term{constant-sheaf}{常值层}。置
\SourcePage{473}
\[
\overline{\mathrm{fr}}_X=\mathrm{fr}_X\times\operatorname{id}_{\bar e},
\]
即 $\mathrm{fr}_X$ 沿 $\bar X\to X$ 的\Term{inverse-image}{逆像}；
\[
f_X=\operatorname{id}_X\times\mathrm{fr}_{\bar e},
\]
即由 $f$ 经结构搬移定义的 $\bar X$ 的\Term{automorphism}{自同构}；
\[
\overline{\mathrm{Fr}}_{\bar F/X}^{*}=\mathrm{Fr}_{F/X}^{*}\boxtimes\operatorname{id}_{\Lambda_{\bar e}},
\]
即 $\mathrm{Fr}_{F/X}^{*}$ 沿 $\bar X\to X$ 的\Term{inverse-image}{逆像}。
可写成
\[
\mathrm{fr}_{\bar X}=\overline{\mathrm{fr}}_X\circ f_X=f_X\circ\overline{\mathrm{fr}}_X,
\qquad \mathrm{Fr}_{\bar F/\bar X}^{*}=\overline{\mathrm{Fr}}_{\bar F/X}^{*}.
\]
命题 3 的推论表明，由二元组 $(\overline{\mathrm{fr}}_X,\overline{\mathrm{Fr}}_{\bar F/X}^{*})$ 定义的 $\mathrm R\Gamma_{\bar X}(\bar F)$ 的\Term{endomorphism}{自同态}，是 $f_X$ 所定义的\Term{automorphism}{自同构}的逆。称 $(\overline{\mathrm{fr}}_X,\overline{\mathrm{Fr}}_{\bar F/X}^{*})$ 为\Term{geometric-frobenius-correspondence}{几何 Frobenius 对应}（它是 $(X,F)$ 上的\Term{frobenius-correspondence}{Frobenius 对应}沿 $\bar X\to X$ 的\Term{inverse-image}{逆像}），并将它所定义的 $\mathrm R\Gamma_{\bar X}(\bar F)$ 的\Term{automorphism}{自同构}记为 $\mathrm{fr}_{\mathrm R\Gamma_{\bar X}(\bar F)}$；由 $f\in\pi_e$ 经结构搬移产生的\Term{arithmetic-frobenius}{算术 Frobenius 作用}是它的逆。

\SGAsection{3}{L 函数}
\textbf{1. 定义。}

设 $p$ 为素数，取素数 $\ell\ne p$ 及 $\mathbb Q_\ell$ 的一个\Term{finite-extension}{有限扩张} $\Omega$。对每个二元组 $(X,F)$，其中 $X$ 是 $e=\operatorname{Spec}(\mathbb F_p)$ 上的\Term{finite-type}{有限型}\Term{scheme}{概形}，$F$ 是 $X$ 上的\Term{constructible}{可构造} $\Omega$-\Term{sheaf}{层}，我们将赋予一个常数项为 $1$ 的\Term{formal-power-series}{形式幂级数} $L_F\in\Omega[[t]]$。
\SourcePage{474}
先考虑 $X=e$ 的特例。设 $\overline{\mathbb F}_p$ 为 $\mathbb F_p$ 的\Term{algebraic-closure}{代数闭包}，$\bar e=\operatorname{Spec}(\overline{\mathbb F}_p)$。\Term{sheaf}{层} $F$ 完全由其\Term{stalk}{茎} $F_{\bar e}$ 决定，后者是带有\Term{galois-group}{Galois 群} $\pi_e=\operatorname{Gal}(\overline{\mathbb F}_p/\mathbb F_p)$ 的\Term{continuous-action}{连续作用}的\Term{finite-dimensional}{有限维} $\Omega$-\Term{vector-space}{向量空间}。由于 $\pi_e$ 由 Frobenius 元素 $f:\lambda\mapsto\lambda^p$（$\lambda\in\overline{\mathbb F}_p$）\Term{topology}{拓扑}生成，给定 $F$ \Term{equivalence}{等价}于给\Term{orientation}{定向}量\Term{space}{空间} $F_{\bar e}$ 以及一个“\Term{topologically-unipotent}{拓扑幺幂}”的\Term{automorphism}{自同构} $f_{F_{\bar e}}$：这就是说，当 $k$ 沿“\Term{multiplication}{乘法}意义”趋于无穷时，其 $k$ 次幂在解析群 $\operatorname{Aut}(F_{\bar e})$ 中趋于恒等。置
\[
L_F(t)=1/\det(1-f_{F_{\bar e}}^{-1}t)\in\Omega[[t]].
\]
现设 $X=\operatorname{Spec}(\mathbb F_q)$，$q=p^d$（整数 $d\geq1$）。设 $\bar x$ 是由 $\mathbb F_q$ 到 $\overline{\mathbb F}_p$ 的一个\Term{embedding}{嵌入}定义的 $X$ 的\Term{geometric-point}{几何点}；\Term{galois-group}{Galois 群} $\pi_x=\operatorname{Gal}(\overline{\mathbb F}_p/\mathbb F_q)$ 可认同于 $\pi_e$ 中由 $f^d$ \Term{topology}{拓扑}\Term{generated-subgroup}{生成的子群}。因此给定\Term{sheaf}{层} $F$ \Term{equivalence}{等价}于给定它的\Term{stalk}{茎} $F_{\bar x}$（\Term{finite-dimensional}{有限维} $\Omega$-\Term{vector-space}{向量空间}），以及其上\Term{topologically-unipotent}{拓扑幺幂}的\Term{automorphism}{自同构} $(f^d)_{F_{\bar x}}$，简记为 $f_{F_{\bar x}}^d$。设 $g$ 为 $X$ 到 $e$ 的唯一\Term{morphism}{态射}。\Term{direct-image}{直像} $g_*(F)$ 同样由其带有 $\pi_e$ \Term{action}{作用}的\Term{stalk}{茎} $g_*(F)_{\bar e}$ 定义；后者按“\Term{induced-module}{诱导模}”计算：
\[
g_*(F)_{\bar e}\simeq F_{\bar x}\otimes_{\pi_x}\pi_e.
\]
置
\[
L_F(t)=L_{g_*(F)}(t)=1/\det(1-f_{g_*(F)_{\bar e}}^{-1}t).
\]
容易计算\Term{determinant}{行列式}，得到此表达式等于
\SourcePage{475}
\[
L_F(t)=1/\det(1-f_{F_{\bar x}}^{-d}t^d).
\]
对一般的 $e$ 上\Term{finite-type}{有限型}\Term{scheme}{概形} $X$ 及其上的\Term{constructible}{可构造} $\Omega$-\Term{sheaf}{层} $F$，定义
\[
L_F=\prod_{x\in X^0}L_{F_x},
\]
其中 $X^0$ 是 $X$ 的\Term{closed-point}{闭点}集。我们将看到此乘积在 $\Omega[[t]]$ 中有意义。具体而言，对每个 $x\in X^0$，取一个位于 $x$ 的 $X$ 的\Term{geometric-point}{几何点} $\bar x$，由\Term{residue-field}{剩余域} $k(x)$ \Term{embedding}{嵌入} $\overline{\mathbb F}_p$ 定义，并置
\[
d(x)=[k(x):\mathbb F_p]\qquad\text{（$x$ 的次数）}.
\]
于是
\[
L_{F_x}(t)=1/\det(1-f_{F_{\bar x}}^{-d(x)}t^{d(x)})
=1+\operatorname{Tr}(f_{F_{\bar x}}^{-d(x)})t^{d(x)}+\cdots.
\]
对每个整数 $n$，满足 $d(x)\leq n$ 的 $x\in X^0$ 只有\Term{finite}{有限}个，所以族 $(L_{F_x})_{x\in X^0}$ 在 $\Omega[[t]]$ 中可乘。因此 $(X,F)$ 的 \Term{l-function}{$L$ 函数}为
\[
L_F(t)=\prod_{x\in X^0}\bigl(1/\det(1-f_{F_{\bar x}}^{-d(x)}t^{d(x)})\bigr).
\]
\Term{l-function}{$L$ 函数}最初的\Term{form}{形式}性质汇集于下一命题。
\begin{SGAresult}{命题}{1}
a)（对 $F$ 的\Term{multiplication}{乘法}性）设 $0\to F'\to F\to F''\to0$ 是 $X$ 上\Term{constructible}{可构造} $\Omega$-\Term{sheaf}{层}的\Term{exact-sequence}{正合列}，
\SourcePage{476}
则 $L_F=L_{F'}L_{F''}$。

b) 设 $Y$ 是 $X$ 的闭子\Term{scheme}{概形}，$U=X-Y$ 是其补\Term{open-subset}{开集}。则对每个\Term{constructible}{可构造} $\Omega$-\Term{sheaf}{层} $F$，有 $L_F=L_{F|U}L_{F|Y}$。

c) 设 $X\to S$ 是 $\mathbb F_p$ 上\Term{finite-type}{有限型}\Term{scheme}{概形}之间的\Term{morphism}{态射}，$F$ 是 $X$ 上\Term{constructible}{可构造} $\Omega$-\Term{sheaf}{层}。则
\[
L_F=\prod_{s\in S^0}L_{F|X_s}.
\]
\end{SGAresult}
\begin{proof}
a) 显然；b) 来自\Term{disjoint-union}{不交并} $X^0=U^0\cup Y^0$；c) 来自 $X^0=\coprod_{s\in S^0}X_s^0$ 及可乘的无穷乘积的\Term{associative-law}{结合律}。上述关于\Term{closed-point}{闭点}的等式成立，是因为 $\mathbb F_p$ 上\Term{finite-type}{有限型}\Term{prescheme}{预概形}之间的\Term{morphism}{态射}把\Term{closed-point}{闭点}仍映为\Term{closed-point}{闭点}。
\end{proof}
\textbf{2. 有理性\Term{theorem}{定理}的陈述及向\Term{curve}{曲线}情形的\Term{reduced}{约化}。}
\begin{SGAresult}{\Term{theorem}{定理}}{}
设 $X$ 是 $\mathbb F_p$ 上\Term{dimension}{维数}为 $n$ 的\Term{finite-type}{有限型}\Term{scheme}{概形}，$g:X\to e$ 为结构\Term{morphism}{态射}，$F$ 是 $X$ 上的\Term{constructible}{可构造} $\Omega$-\Term{sheaf}{层}。则
\[
L_F=\prod_{i=0}^{2n}(L_{\mathrm R^ig_!(F)})^{(-1)^i}.
\tag{1}
\]
\end{SGAresult}
这里 $\mathrm R^ig_!(F)$ 表示 $\ell$ 进\Term{cohomology}{上同调}意义下的\Term{proper-support-direct-image}{固有支撑高阶直像}。置
\[
H_!^i(\bar X,\bar F)=\mathrm R^i\bar g_!(\bar F)=\mathrm R^ig_!(F)_{\bar e}
\]
（\Term{base-change-theorem}{基变换定理}），其中
\SourcePage{477}
$\bar X$、$\bar F$、$\bar g$ 分别由 $X$、$F$、$g$ 经\Term{base-change}{基变换} $\bar e\to e$ 得到（$\bar e$ 为 $\mathbb F_p$ 的\Term{algebraic-closure}{代数闭包}的谱）。用此记号，由第 1 小节有
\[
L_{\mathrm R^ig_!(F)}(t)
=1/\det(1-f_{H_!^i(\bar X,\bar F)}^{-1}t)
=L_{f_{H_!^i(\bar X,\bar F)}^{-1}}(t),
\]
其中，遵循前面某讲建立的\Term{form}{形式}体系，对\Term{finite-dimensional-vector-space}{有限维向量空间}的每个\Term{endomorphism}{自同态} $u$，置 $L_u(t)=1/\det(1-ut)$。

更一般地，考虑 $\mathbb F_p$ 上\Term{finite-type}{有限型}\Term{scheme}{概形}之间的\Term{morphism}{态射} $u:X\to Y$ 以及 $X$ 上的\Term{constructible}{可构造} $\Omega$-\Term{sheaf}{层} $F$。以 $(P_{F,u})$ 表示性质
\[
(P_{F,u}):\qquad L_F=\prod_i(L_{\mathrm R^iu_!(F)})^{(-1)^i}.
\]
\Term{theorem}{定理}对 $(X,F)$ 的陈述正是 $(P_{F,g})$，其中 $g:X\to e$ 为 $X$ 的结构\Term{morphism}{态射}。下面借助若干\Term{lemma}{引理}，将它\Term{reduced}{约化}到 $X$ 为\Term{projective-line}{射影直线}的情形。
\begin{SGAresult}{\Term{lemma}{引理}}{1}
设 $u:X\to Y$ 是 $\mathbb F_p$ 上\Term{finite-type}{有限型}\Term{scheme}{概形}之间的\Term{morphism}{态射}，$F$ 是\Term{constructible}{可构造} $\Omega$-\Term{sheaf}{层}。若对所有 $y\in Y^0$，$(P_{F|X_y,u_y})$ 成立，则 $(P_{F,u})$ 成立。
\end{SGAresult}
\begin{proof}
确实，
\[
\begin{aligned}
L_F&=\prod_{y\in Y^0}L_{F|X_y}&&\text{（第 1 小节命题 1 c)）}\\
&=\prod_{y\in Y^0}\prod_i(L_{\mathrm R^iu_{y!}(F|X_y)})^{(-1)^i}
&&\text{（使用 $(P_{F|X_y,u_y})$）}.
\end{aligned}
\]
而 $\mathrm R^iu_{y!}(F|X_y)\simeq\mathrm R^iu_!(F)_y$（\Term{base-change-theorem}{基变换定理}），故
\SourcePage{478}
\[
L_F=\prod_i\left(\prod_{y\in Y^0}L_{\mathrm R^iu_!(F)_y}\right)^{(-1)^i}
=\prod_i(L_{\mathrm R^iu_!(F)})^{(-1)^i},
\]
最后一步使用 $Y$ 上\Term{l-function}{$L$ 函数}的定义。
\end{proof}
\begin{SGAresult}{推论}{}
设 $u:X\to Y$ 是 $\mathbb F_p$ 上\Term{finite-type}{有限型}\Term{scheme}{概形}之间的\Term{finite-morphism}{有限态射}，$F$ 是 $X$ 上\Term{constructible}{可构造} $\Omega$-\Term{sheaf}{层}，则性质 $(P_{F,u})$ 成立。
\end{SGAresult}
\begin{proof}
由\Term{lemma}{引理}\Term{reducible}{可约}化到 $Y$ 为单点，即 $Y=\operatorname{Spec}(\mathbb F_q)$。此时 $X$ 为\Term{finite}{有限}的，并且 $u_*(F)\simeq\bigoplus_{x\in X}u_*(F_x)$。故为证明 $(P_{F,u})$ 的公式 $L_F=L_{u_*(F)}$，可进一步\Term{reduced}{约化}到 $X$ 本身为单点：$X=\operatorname{Spec}(\mathbb F_{q'})$。设 $g:X\to e$、$h:Y\to e$ 为结构\Term{morphism}{态射}。由于 $g=h\circ u$，按定义（第 1 小节）有
\[
L_F=L_{g_*(F)}=L_{h_*u_*(F)}=L_{u_*(F)}.
\]
\end{proof}
\begin{SGAresult}{\Term{lemma}{引理}}{2}
设 $u:X\to Y$、$v:Y\to Z$ 是 $\mathbb F_p$ 上\Term{finite-type}{有限型}\Term{scheme}{概形}之间的\Term{morphism}{态射}，$F$ 是 $X$ 上\Term{constructible}{可构造} $\Omega$-\Term{sheaf}{层}。若 $(P_{F,u})$ 以及对所有 $i$ 的 $(P_{\mathrm R^iu_!(F),v})$ 均成立，则 $(P_{F,v\circ u})$ 成立。
\end{SGAresult}
\begin{proof}
确实，由 $(P_{F,u})$ 与 $(P_{\mathrm R^iu_!(F),v})$，有
\[
\begin{aligned}
L_F&=\prod_i(L_{\mathrm R^iu_!(F)})^{(-1)^i}\\
&=\prod_i\left(\prod_j(L_{\mathrm R^jv_!(\mathrm R^iu_!(F))})^{(-1)^j}\right)^{(-1)^i}\\
&=\prod_{i,j}(L_{\mathrm R^jv_!(\mathrm R^iu_!(F))})^{(-1)^{i+j}}.
\end{aligned}
\]
Leray \Term{spectral-sequence}{谱序列} $\mathrm R^jv_!(\mathrm R^iu_!(F))\Rightarrow\mathrm R^k(v\circ u)_!(F)$ 及 $L_F$ 对 $F$ 的\Term{multiplication}{乘法}性（第 1 小节命题 1 a)）最终给出
\SourcePage{479}
\[
L_F=\prod_k(L_{\mathrm R^k(v\circ u)_!(F)})^{(-1)^k},
\]
即所欲证。
\end{proof}
\begin{SGAresult}{\Term{lemma}{引理}}{3}
设 $X$ 是 $\mathbb F_p$ 上\Term{finite-type}{有限型}\Term{scheme}{概形}，$g:X\to e$ 为结构\Term{morphism}{态射}。

a) 考虑 $X$ 上的\Term{constructible}{可构造} $\Omega$-\Term{sheaf}{层} $F$、$X$ 的闭子\Term{scheme}{概形} $Y$ 及其补\Term{open-subset}{开集} $U=X-Y$。若\Term{theorem}{定理}对 $(X,F)$、$(Y,F|Y)$、$(U,F|U)$ 中的两个成立，则对第三个也成立。

b) 考虑 $X$ 上\Term{constructible}{可构造} $\Omega$-\Term{sheaf}{层}的\Term{exact-sequence}{正合列} $0\to F'\to F\to F''\to0$。若\Term{theorem}{定理}对 $(X,F')$、$(X,F)$、$(X,F'')$ 中的两个成立，则对第三个也成立。
\end{SGAresult}
\begin{proof}
我们只会用到 a)，但出于\Term{symmetric}{对称}性也给出 b)。a)、b) 的证明类似，例如证明 a)。\Term{exact-sequence}{正合列}
\[
\cdots\longrightarrow\mathrm R^ig_{U!}(F|U)\longrightarrow\mathrm R^ig_!(F)
\longrightarrow\mathrm R^ig_{Y!}(F|Y)\longrightarrow\mathrm R^{i+1}g_{U!}(F|U)\longrightarrow\cdots
\]
与命题 1 a)（第 1 小节）给出
\[
\prod_i(L_{\mathrm R^ig_!(F)})^{(-1)^i}
=\prod_i\bigl(L_{\mathrm R^ig_{U!}(F|U)}L_{\mathrm R^ig_{Y!}(F|Y)}\bigr)^{(-1)^i},
\]
再由命题 1 b)（第 1 小节）即得结论。
\end{proof}
\Term{lemma}{引理} 1、2、3 的论证可以概括为：根据第 1 小节命题 1，\Term{l-function}{$L$ 函数} $L_F$ 具有与\Term{proper-support-cohomology}{固有支撑上同调}类似的\Term{form}{形式}性质。
\SourcePage{480}
现利用上述\Term{lemma}{引理}\Term{reduced}{约化}到 $X$ 为\Term{projective-line}{射影直线}的情形。

\textbf{第一次\Term{reduced}{约化}。}由\Term{noetherian}{诺特}归纳，可将 $X$ \Term{partition}{分割}为\Term{affine}{仿射}\Term{scheme}{概形} $(U_j)_{1\leq j\leq r}$，使 $U_{j+1}$ 为 $X-\bigcup_{i=1}^jU_i$ 中的\Term{open-subset}{开集}（$j=1,\ldots,r-1$）。假设对每个二元组 $(U,G)$，其中 $U$ 是 $\mathbb F_p$ 上\Term{finite-type}{有限型}\Term{affine}{仿射}\Term{scheme}{概形}、$G$ 是其上的\Term{constructible}{可构造} $\Omega$-\Term{sheaf}{层}，\Term{theorem}{定理}均已成立。则\Term{theorem}{定理}对 $(U_j,F|U_j)$（$j=1,\ldots,r$）成立。由\Term{lemma}{引理} 3 a)，对 $j$ 作\Term{descent}{下降}归纳，得到它对
\[
\left(\bigcup_{i\geq j}U_i,\,F\bigm|\bigcup_{i\geq j}U_i\right)
\]
成立，因此对 $(X,F)$ 成立，因为 $X=\bigcup_{i\geq1}U_i$。所以只须证明 $X$ \Term{affine}{仿射}的情形。

\textbf{第二次\Term{reduced}{约化}。}设 $X$ \Term{affine}{仿射}且\Term{dimension}{维数}为 $n$。存在从 $X$ 到 $n$ \Term{dimension}{维数}值\Term{space}{空间}的\Term{finite-morphism}{有限态射} $u:X\to\mathbf E^n$。对 $n$ 归纳。$n=0$ 时，$X$ 在 $e$ 上\Term{finite}{有限}，\Term{theorem}{定理}由\Term{lemma}{引理} 1 的推论成立。若 $n\geq1$，假设\Term{theorem}{定理}对\Term{dimension}{维数}为 $n-1$ 的\Term{affine}{仿射}\Term{scheme}{概形}成立。取\Term{finite-morphism}{有限态射} $u:X\to\mathbf E^n$，及由\Term{isomorphism}{同构} $\mathbf E^n\simeq\mathbf E^1\times\mathbf E^{n-1}$ 定义的\Term{projection}{投影} $v:\mathbf E^n\to\mathbf E^1$。复合得到 $w=v\circ u:X\to\mathbf E^1$，其\Term{fiber}{纤维}是\Term{dimension}{维数}为 $n-1$ 的\Term{affine}{仿射}\Term{scheme}{概形}。由归纳假设及\Term{lemma}{引理} 1，性质 $(P_{F,w})$ 成立。借助\Term{lemma}{引理} 2，可见只需对二元组
\SourcePage{481}
$(\mathbf E^1,\mathrm R^iw_!(F))$ 证明\Term{theorem}{定理}。这样\Term{reduced}{约化}到了 $X=\mathbf E^1$ 为\Term{affine-line}{仿射直线}的情形。

\textbf{第三次\Term{reduced}{约化}。}设 $X=\mathbf E^1$，$i:\mathbf E^1\to\mathbf P^1$ 是\Term{affine-line}{仿射直线}到\Term{projective-line}{射影直线} $\mathbf P^1$ 的\Term{immersion}{浸入}，$Z=\mathbf P^1-i(\mathbf E^1)$ 是由 $\mathbf P^1$ 的无穷远点组成的\Term{reduced}{约化}闭子\Term{scheme}{概形}。由于\Term{theorem}{定理}对 $(Z,i_!(F)|Z)$ 成立，根据\Term{lemma}{引理} 3 a)，对 $(\mathbf E^1,F)$ 证明它\Term{equivalence}{等价}于对 $(\mathbf P^1,i_!(F))$ 证明它（因为 $i$ 是\Term{open-immersion}{开浸入}，有 $F=i_!(F)|\mathbf E^1$）。

\textbf{3. \Term{theorem}{定理}的证明。}

经过上述\Term{reduced}{约化}，只须对 $X$ 为\Term{projective-line}{射影直线}的情形证明\Term{theorem}{定理}。先将\Term{theorem}{定理}中的公式（1）改写成关于迹的 Lefschetz 型公式；再在 $X$ 为\Term{smooth}{光滑}\Term{projective}{射影}\Term{curve}{曲线}的特例下证明后一公式。

重写公式（1）：
\[
L_F=\prod_i(L_{\mathrm R^ig_!(F)})^{(-1)^i}
=\frac{P_F^1\cdots P_F^{2n-1}}{P_F^0\cdots P_F^{2n}},
\qquad P_F^i(t)=\det(1-f_{H_!^i(\bar X,\bar F)}^{-1}t).
\]
左边的对数为
\[
\begin{aligned}
\sum_{x\in X^0}\sum_{n=1}^{\infty}\operatorname{Tr}(f_{F_{\bar x}}^{-nd(x)})\frac{t^{nd(x)}}n
=\sum_{n=1}^{\infty}\left(\sum_{d(x)\mid n}d(x)\operatorname{Tr}(f_{F_{\bar x}}^{-n})\right)\frac{t^n}n,
\end{aligned}
\]
而右边的对数为
\SourcePage{482}
\[
\sum_{n\geq1}\left(\sum_i(-1)^i\operatorname{Tr}(f_{H_!^i(\bar X,\bar F)}^{-n})\right)\frac{t^n}n.
\]
比较 $t^n/n$ 的系数，可见须证明：对所有 $n\geq1$，
\[
\sum_{d(x)\mid n}d(x)\operatorname{Tr}(f_{F_{\bar x}}^{-n})
=\sum_i(-1)^i\operatorname{Tr}(f_{H_!^i(\bar X,\bar F)}^{-n}).
\tag{2}
\]
这些计算合法，是因为域 $\Omega$ 的\Term{characteristic}{特征}为零。

再变换公式（2），假设 $X$ \Term{proper-scheme}{固有}，从而可将\Term{proper-support-cohomology}{固有支撑上同调}换成普通\Term{cohomology}{上同调}。回忆，$X$ 上 $\mathbb Q_\ell$-\Term{sheaf}{层}的\Term{category}{范畴}是 $\ell$ 进\Term{sheaf}{层}\Term{category}{范畴}的一个分式\Term{category}{范畴}。从现在起，略滥用记号，用 $F=(F_\nu)_{\nu\in\mathbb N}$ 表示 $X$ 上与 $F$ 所在的 $\mathbb Q_\ell$-\Term{sheaf}{层}相对应的 $\ell$ 进\Term{sheaf}{层}\footnote{为简化书写，下文的证明假设 $\Omega=\mathbb Q_\ell$。一般情形只须将 $\mathbb Z_\ell$ 换成它在 $\Omega$ 中的\Term{integral-closure}{整闭包}。（细节留给读者。）}。于是 $\bar F=(\bar F_\nu)$，而 Frobenius \Term{morphism}{态射} $\mathrm{Fr}_{F_\nu/X}^{*}$ 定义 $\ell$ 进\Term{sheaf}{层}的\Term{morphism}{态射}
\[
\overline{\mathrm{Fr}}_{\bar F/X}^{*}:\overline{\mathrm{fr}}_X^*(\bar F)\longrightarrow\bar F,
\]
从而得到 $(\bar X,\bar F)$ 上的 $\ell$ 进\Term{correspondence}{对应}。对每个整数 $\nu$，我们知道由 $(\overline{\mathrm{fr}}_X,\overline{\mathrm{Fr}}_{\bar F_\nu/X}^{*})$ 定义的\Term{endomorphism}{自同态} $\mathrm{fr}_{\mathrm R\Gamma_{\bar X}(\bar F_\nu)}$ 是 $f_{\mathrm R\Gamma_{\bar X}(\bar F_\nu)}$ 的逆（第 2 节第 3 小节命题 3 的推论）。利用 $\ell$ 进\Term{cohomology}{上同调}的定义，可得由 $(\overline{\mathrm{fr}}_X,\overline{\mathrm{Fr}}_{\bar F/X}^{*})$ 定义的 $H^i(\bar X,\bar F)$ 的\Term{endomorphism}{自同态} $\mathrm{fr}_{H^i(\bar X,\bar F)}$ 是 $f_{H^i(\bar X,\bar F)}$ 的逆。因此，对所有 $n\geq1$ 及所有 $i$，
\[
f_{H^i(\bar X,\bar F)}^{-n}=\mathrm{fr}_{H^i(\bar X,\bar F)}^n.
\]
将此结果用于 $X=x=\operatorname{Spec}(\mathbb F_q)$（$q=p^d$）\Term{reduced}{约化}为单点且\Term{degree}{次数} $d$ 整除 $n$ 的特例，得到
\SourcePage{483}
\[
f_{F_{\bar x}}^{-n}=\mathrm{fr}_{F_{\bar x}}^n
=(\overline{\mathrm{Fr}}_{\bar F/X}^{*})_{\bar x}^{\,n},
\]
其中 $\bar x$ 是由 $\mathbb F_q$ 到 $\overline{\mathbb F}_p$ 的一个\Term{embedding}{嵌入}定义的\Term{geometric-point}{几何点}。

另一方面，$\overline{\mathrm{fr}}_X$ 与 $f_X$ 在 $\bar X$ 的点上互为逆\Term{action}{作用}。因此\Term{projection}{投影} $\bar X\to X$ 在 $X$ 的\Term{degree}{次数}为 $d$ 的\Term{closed-point}{闭点}集与 $\overline{\mathrm{fr}}_X$ 在 $\bar X$ 中\Term{action}{作用}的 $d$ 元\Term{orbit}{轨道}集之间建立\Term{bijection}{双射}。注意，$\bar X$ 中被 $\overline{\mathrm{fr}}_X^{\,n}$ 固定的点集 $\bar X^{\overline{\mathrm{fr}}_X^{\,n}}$，正是 $\overline{\mathrm{fr}}_X$ 的\Term{cardinal}{基数}整除 $n$ 的那些\Term{orbit}{轨道}的并，故可将（2）改写成“几何”\Term{form}{形式}：
\[
\sum_{\bar x\in\bar X^{\overline{\mathrm{fr}}_X^{\,n}}}
\operatorname{Tr}(\mathrm{fr}_{\bar F_{\bar x}}^n)
=\sum_i(-1)^i\operatorname{Tr}(\mathrm{fr}_{H^i(\bar X,\bar F)}^n).
\tag{2'}
\]
这就是\Term{scheme}{概形} $\bar X$、\Term{sheaf}{层} $\bar F$ 与\Term{correspondence}{对应} $(\overline{\mathrm{fr}}_X^n,\overline{\mathrm{Fr}}_{\bar F/X}^{*\,n})$ 的 Lefschetz 型公式。

为简化记号，从现在起，将 $\bar X$、$\bar F$、$\overline{\mathrm{fr}}_X^n$ 及 $(\overline{\mathrm{Fr}}_{\bar F/X}^{*})^n$ 分别写为 $X$、$F$、$h$ 及 $h_F$（原来的 $X$、$F$ 已不再出现，不会混淆）。假设 $X$ 是 $k=\overline{\mathbb F}_p$ 上的\Term{smooth}{光滑}\Term{projective}{射影}\Term{curve}{曲线}。$h$ 在 $X$ 中的\Term{fixed-point}{不动点}的\Term{multiplicity}{重数}均为 $1$；确实，在\Term{fixed-point}{不动点}的\Term{local-ring}{局部环}的\Term{completion}{完备化}上，$h$ 的\Term{action}{作用}为 $\tau\mapsto\tau^{p^n}$，其中 $\tau$ 是一个\Term{uniformizer}{素元}。\Term{theorem}{定理}将由下述结果\Term{pushout}{推出}。
\begin{SGAresult}{命题}{2}
设 $X$ 是\Term{characteristic}{特征}为 $p$ 的\Term{algebraically-closed-field}{代数闭域} $k$ 上的\Term{smooth}{光滑}\Term{projective}{射影}\Term{curve}{曲线}，$F$ 是 $X_{\text{ét}}$ 上的
\SourcePage{484}
\Term{constructible}{可构造} $\ell$ 进\Term{sheaf}{层}（$\ell$ 为不等于 $p$ 的素数），$h$ 是 $X$ 的\Term{endomorphism}{自同态}，其所有\Term{fixed-point}{不动点}的\Term{multiplicity}{重数}均为 $1$，且 $h_F:h^*(F)\to F$ 为\Term{sheaf}{层}\Term{morphism}{态射}。以下公式成立，其中 $h_{H^i(X,F)}$ 表示由 $(h,h_F)$ 定义的 $H^i(X,F)$ 的\Term{endomorphism}{自同态}：
\[
\sum_{x\in X^h}\operatorname{Tr}(h_{F_x})
=\sum_i(-1)^i\operatorname{Tr}(h_{H^i(X,F)}).
\tag{2''}
\]
\end{SGAresult}
\begin{proof}
本讨论班已经对\Term{curve}{曲线} $X$ 上被 $\ell^{\nu+1}$ 消去的\Term{constructible}{可构造}\Term{torsion-sheaf}{挠层} $F_\nu$，连同\Term{morphism}{态射} $h_{F_\nu}:h^*(F_\nu)\to F_\nu$，建立过类似的 Lefschetz 公式：
\[
\sum_{x\in X^h}\operatorname{Tr}(h_{(F_\nu)_x})
=\operatorname{Tr}^{*}h_{\mathrm R\Gamma_X(F_\nu)},
\tag{3}
\]
等式两边均属于 $\mathbb Z/\ell^{\nu+1}\mathbb Z$。

将命题中的 $\ell$ 进\Term{sheaf}{层}写成 $F=(F_\nu)_{\nu\in\mathbb N}$，其中 $F_\nu$ 被 $\ell^{\nu+1}$ 消去。对每个 $\nu$ 都可写出公式 $(3)_\nu$；这些公式左边随 $\nu$ 变化，正是某个 $\ell$ 进整数的各\Term{set-component}{分量}，而这个整数恰好是待证公式（2''）的左边。还须证明：各公式 $(3)_\nu$ 的右边随 $\nu$ 变化，定义的 $\ell$ 进整数等于 $\sum_i(-1)^i\operatorname{Tr}(h_{H^i(X,F)})$，其中
\[
H^i(X,F)=\varprojlim_\nu H^i(\mathrm R\Gamma_X(F_\nu)).
\]
这一点由最后要证明的若干纯\Term{algebra}{代数}\Term{lemma}{引理}\Term{pushout}{推出}。

考虑\Term{perfect-complex}{完美复形}的\Term{projective-system}{投射系} $(K_\nu)_{\nu\in\mathbb N}$，其中 $K_\nu\in D_{\mathrm{parf}}(\mathbb Z/\ell^{\nu+1}\mathbb Z)$，并且对每个 $\nu$，转移\Term{morphism}{态射}
\SourcePage{485}
$K_{\nu+1}\to K_\nu$ 在\Term{derived-category}{导出范畴}意义下\Term{isomorphism}{同构}于\Term{canonical}{典范}\Term{morphism}{态射} $K_{\nu+1}\to K_{\nu+1}/\ell^{\nu+1}K_{\nu+1}$。假设每个\Term{complex}{复形} $K_\nu$ 带有\Term{endomorphism}{自同态} $h_\nu\in\operatorname{Fl}(D_{\mathrm{parf}}(\mathbb Z/\ell^{\nu+1}\mathbb Z))$，使 $(h_\nu)_{\nu\in\mathbb N}$ 与转移\Term{morphism}{态射}相容，从而定义了\Term{projective-system}{投射系} $(K_\nu)$ 的\Term{endomorphism}{自同态}。在证明命题 2 的应用中，取 $K_\nu=\mathrm R\Gamma_X(F_\nu)$、$h_\nu=h_{\mathrm R\Gamma_X(F_\nu)}$；关于转移\Term{morphism}{态射}的条件，由 $\ell$ 进\Term{projective-system}{投射系} $(F_\nu)$ 的转移\Term{morphism}{态射}所满足的类似条件，通过“\Term{kunneth-formula}{Künneth 公式}”得到。

我们要找一个\Term{projective-system}{投射系} $(K'_\nu)_{\nu\in\mathbb N}$，其中 $K'_\nu$ 是由\Term{finite-type}{有限型}\Term{free-module}{自由} $\mathbb Z/\ell^{\nu+1}\mathbb Z$-\Term{module}{模}组成的\Term{complex}{复形}，在一个与 $\nu$ 无关的\Term{degree}{次数}区间 $[a,b]$ 外为零，并带有\Term{complex}{复形}\Term{endomorphism}{自同态} $h'_\nu:K'_\nu\to K'_\nu$，使下列条件成立：

1) 对每个 $\nu\in\mathbb N$，转移\Term{morphism}{态射} $K'_{\nu+1}\to K'_\nu$ 作为\Term{complex}{复形}\Term{morphism}{态射}，\Term{isomorphism}{同构}于 $K'_{\nu+1}\to K'_{\nu+1}/\ell^{\nu+1}K'_{\nu+1}$。

2) $(h'_\nu)_{\nu\in\mathbb N}$ 与转移\Term{morphism}{态射}相容。

3) 对每个 $\nu$，$(K'_\nu,h'_\nu)$ 与 $(K_\nu,h_\nu)$ 在 $D_{\mathrm{parf}}(\mathbb Z/\ell^{\nu+1}\mathbb Z)$ 中\Term{isomorphism}{同构}，且这些\Term{isomorphism}{同构}可选择得与 $(K'_\nu)$ 及 $(K_\nu)$ 的转移\Term{morphism}{态射}相容。

假设系 $(K'_\nu)$ 与 $(h'_\nu)$ 已构造好。对每个 $\nu\in\mathbb N$，可写
\[
\operatorname{Tr}^{*}h_\nu=\operatorname{Tr}^{*}h'_\nu
=\sum_i(-1)^i\operatorname{Tr}h_\nu'{}^i.
\]
置 $K'=\varprojlim_\nu K'_\nu$，$h'=\varprojlim_\nu h'_\nu$（这是 $K'$ 的\Term{endomorphism}{自同态}）。可见
\SourcePage{486}
$(\operatorname{Tr}^{*}h_\nu)_\nu$ 是 $\ell$ 进整数 $\sum_i(-1)^i\operatorname{Tr}h'{}^i$。此外，\Term{projective-system}{投射系} $(K'_\nu)$ 满足 \Term{ml-condition}{Mittag--Leffler 条件}，对每个 $i$，系 $(H^i(K'_\nu))_\nu$ 亦然；因而
\[
H^i(K')\xrightarrow{\sim}\varprojlim_\nu H^i(K'_\nu)
\qquad\text{对所有 }i
\]
（EGA $0_{\mathrm{III}}$ 13.2.3）。在所需的情形 $K_\nu=\mathrm R\Gamma_X(F_\nu)$、$h_\nu=h_{\mathrm R\Gamma_X(F_\nu)}$ 下，$H^i(K')$ 可认同于 $H^i(X,F)$，而 $h'$ 所定义的 $H^i(K')$ 的\Term{endomorphism}{自同态}可认同于 $h_{H^i(X,F)}$。所以
\[
\sum_i(-1)^i\operatorname{Tr}h'{}^i
=\sum_i(-1)^i\operatorname{Tr}h_{H^i(X,F)},
\]
命题中的公式（2''）因此得证。

系 $(K'_\nu)$ 及 $(h'_\nu)$ 的存在性依赖于两个\Term{lemma}{引理}。在这些\Term{lemma}{引理}中，$A$ 为\Term{commutative-ring}{交换环}，$I$ 为 $A$ 的\Term{ideal}{理想}，$A_0=A/I$，并考虑\Term{functor}{函子} $T:L\mapsto L\otimes_AA_0$（$L$ 为右\Term{bounded}{有界} $A$-\Term{module}{模}\Term{complex}{复形}）。
\end{proof}
\begin{SGAresult}{\Term{lemma}{引理}}{1}
假设\Term{ideal}{理想} $I$ \Term{nilpotent}{幂零}。设 $L\in D^-(A)$ 是各次均\Term{flat}{平坦}的\Term{complex}{复形}，$M\in D^-(A_0)$ 是各次均\Term{projective-module}{投射}的\Term{complex}{复形}，并有 $D^-(A_0)$ 中的\Term{isomorphism}{同构} $u:T(L)\xrightarrow{\sim}M$。则存在右\Term{bounded}{有界}、各次均\Term{projective-module}{投射}的 $A$-\Term{module}{模}\Term{complex}{复形} $L'$，$D^-(A)$ 中的\Term{isomorphism}{同构} $v:L\xrightarrow{\sim}L'$，以及\Term{complex}{复形}\Term{isomorphism}{同构} $w:T(L')\xrightarrow{\sim}M$，使在 $D^-(A_0)$ 中 $u=w\circ T(v)$。
\end{SGAresult}
\textbf{注记。}若 $M$ 各次均\Term{free-module}{自由}（相应地\Term{finite-type}{有限型}），则 $L'$ 各次也均\Term{free-module}{自由}（相应地\Term{finite-type}{有限型}）。
\SourcePage{487}
确实，若 $N$ 是\Term{flat}{平坦} $A$-\Term{module}{模}，而 $T(N)$ \Term{free-module}{自由}（相应地\Term{finite-type}{有限型}），则容易证明 $N$ 也\Term{free-module}{自由}（相应地\Term{finite-type}{有限型}）（Bourbaki，《\Term{commutative-algebra}{交换代数}》第 II 章，第 3 节，命题 5）。由此得到一个推论：若 $L\in\operatorname{Ob}D^-(A)$，使得 $L\mathbin{\overset{\mathrm L}{\otimes}}_AA_0$ 为\Term{pseudo-coherent}{伪凝聚}（相应地\Term{perfect-complex}{完美}）的 $A_0$-\Term{module}{模}\Term{complex}{复形}（即在 $D^-(A)$ 中\Term{isomorphism}{同构}于一个\Term{bounded-above}{上有界}、相应地\Term{bounded}{有界}，且各次为\Term{finite-projective-module}{有限型投射模}的\Term{complex}{复形}），则 $L$ 是\Term{pseudo-coherent}{伪凝聚}（相应地\Term{perfect-complex}{完美}）的 $A$-\Term{module}{模}\Term{complex}{复形}。
\begin{SGAresult}{\Term{lemma}{引理}}{2}
设 $L$ 是右\Term{bounded}{有界}、各次均\Term{projective-module}{投射}的 $A$-\Term{module}{模}\Term{complex}{复形}，带有\Term{endomorphism}{自同态} $\varphi$；设 $\varphi_0$ 是 $M=T(L)$ 的\Term{endomorphism}{自同态}，在\Term{derived-category}{导出范畴} $D^-(A_0)$ 意义下与 $T(\varphi)$ 相同。则存在 $L$ 的\Term{endomorphism}{自同态} $\varphi'$，在 $D^-(A)$ 意义下与 $\varphi$ 相同，且 $T(\varphi')=\varphi_0$。
\end{SGAresult}
在证明这两个\Term{lemma}{引理}之前，先说明如何用它们证明系 $(K'_\nu)$ 与 $(h'_\nu)$ 的存在性。对 $\nu$ 归纳，假设 $K'_0,h'_0,\ldots,K'_{\nu-1},h'_{\nu-1}$ 已构造好。将\Term{lemma}{引理}用于
\[
A=\mathbb Z/\ell^{\nu+1}\mathbb Z,\qquad
I=\ell^\nu\mathbb Z/\ell^{\nu+1}\mathbb Z\quad\text{（故 $I^2=0$）},\qquad
A_0=\mathbb Z/\ell^\nu\mathbb Z.
\]
先在第一个\Term{lemma}{引理}中取 $L=K_\nu$、$M=K'_{\nu-1}$。结合注记 2，得到 $K'_\nu=L'$，转移\Term{morphism}{态射} $K'_\nu\to K'_{\nu-1}$ 由 $w$ 给出，并满足前述条件 1)。$D^-(A)$ 中的\Term{isomorphism}{同构} $v$ 可将 $h_\nu$ 搬移为 $D^-(A)$ 意义下 $K'_\nu$ 的\Term{endomorphism}{自同态}。可假设后者来自\Term{complex}{复形}\Term{endomorphism}{自同态} $h''_\nu:K'_\nu\to K'_\nu$。再将第二个\Term{lemma}{引理}用于 $L=K'_\nu$、
\SourcePage{488}
$\varphi=h''_\nu$、$\varphi_0=h'_{\nu-1}$，便得到满足条件 2) 的 $h'_\nu:K'_\nu\to K'_\nu$，同时条件 3) 也成立。注意，在\Term{lemma}{引理} 1 的记号下，若某个 $i$ 满足 $M^i=0$，则由 \Term{nakayama-lemma}{Nakayama 引理}可得 $L'{}^i=0$。因此所有\Term{complex}{复形} $K'_\nu$ 确实在同一个\Term{degree}{次数}区间以外为零。

为证明\Term{lemma}{引理} 1、2，我们需要下述结果。
\begin{SGAresult}{\Term{lemma}{引理}}{3}
(i) 设 $P$、$Q$ 为 $A$-\Term{module}{模}。若 $P$ \Term{projective-module}{投射}，则对每个 $A$-\Term{linear}{线性}映射 $f_0:T(P)\to T(Q)$，存在 $A$-\Term{linear}{线性}映射 $f:P\to Q$ 使 $T(f)=f_0$。

(ii) 假设\Term{ideal}{理想} $I$ \Term{nilpotent}{幂零}。对每个\Term{projective-module}{投射} $A_0$-\Term{module}{模} $P_0$，存在\Term{projective-module}{投射} $A$-\Term{module}{模} $P$ 使 $T(P)\simeq P_0$。此外，若 $P'$ 是满足 $T(P')\simeq P_0$ 的\Term{flat}{平坦} $A$-\Term{module}{模}，则 $P'\simeq P$。
\end{SGAresult}
\begin{proof}
(i) 显然：若 $u_P$（相应地 $u_Q$）表示从 $P$ 到 $T(P)$（相应地从 $Q$ 到 $T(Q)$）的\Term{canonical-map}{典范映射}，则复合 $f_0\circ u_P$ 是从 $P$ 到 $T(Q)$ 的 $A$-\Term{linear}{线性}映射，且 $u_Q$ \Term{surjective}{满射}。由于 $P$ \Term{projective-module}{投射}，存在经过 $u_Q$ 的\Term{resolution}{分解} $f_0\circ u_P=u_Q\circ f$，其中 $f:P\to Q$ 是 $A$-\Term{linear}{线性}映射；于是 $T(f)=f_0$。

为证明 (ii)，先注意 $P_0$ \Term{free-module}{自由}的情形显然成立（无需对\Term{ideal}{理想} $I$ 作任何假设）。为转到 $P_0$ \Term{projective-module}{投射}的一般情形，取一个\Term{free-module}{自由} $A_0$-\Term{module}{模}
\SourcePage{489}
$L_0$，以 $P_0$ 为\Term{direct-summand}{直和因子}，再取\Term{free-module}{自由} $A$-\Term{module}{模} $L$，使 $T(L)\simeq L_0$。$P_0$ 可认同于\Term{projection-operator}{投影算子} $e_0:L_0\to L_0$ 的像。只须说明：存在\Term{lifting}{提升} $e_0$ 的\Term{projection-operator}{投影算子} $e:L\to L$。这可由 Bourbaki《\Term{commutative-algebra}{交换代数}》第 III 章，第 4 节，第 6 小节，\Term{lemma}{引理} 2 得到：将该\Term{lemma}{引理}用于 $\operatorname{End}(L)$ 中由\Term{lifting}{提升} $e_0$ 的某个\Term{endomorphism}{自同态}生成的 $A$-\Term{subalgebra}{子代数} $B$（该\Term{lifting}{提升}的存在性由 (i) 保证），以及\Term{canonical}{典范}\Term{homomorphism}{同态} $\operatorname{End}(L)\to\operatorname{End}(L_0)$ 的核在 $B$ 上截出的\Term{nilpotent}{幂零}\Term{ideal}{理想}。这样证明了 (ii) 中 $P$ 的存在性。

再设 $P'$ 如 (ii) 所述。由 (i)，复合\Term{isomorphism}{同构} $T(P)\xrightarrow{\sim}P_0\xrightarrow{\sim}T(P')$ 来自某个\Term{homomorphism}{同态} $f:P\to P'$。由于 $P$、$P'$ 在 $A$ 上\Term{flat}{平坦}，且 $T(f)$ 为\Term{isomorphism}{同构}，可得 $\operatorname{Gr}_I(f)$ 为\Term{isomorphism}{同构}。因 $I$ \Term{nilpotent}{幂零}，$f$ 本身是\Term{isomorphism}{同构}，故 (ii) 成立。
\end{proof}
\begin{SGAresult}{推论}{}
假设\Term{ideal}{理想} $I$ \Term{nilpotent}{幂零}。对每个由\Term{projective-module}{投射} $A_0$-\Term{module}{模}组成的右\Term{bounded}{有界}\Term{acyclic-complex}{无上同调复形} $N_0$，存在由\Term{projective-module}{投射} $A$-\Term{module}{模}组成的右\Term{bounded}{有界}\Term{acyclic-complex}{无上同调复形} $N$，使 $T(N)\simeq N_0$。
\end{SGAresult}
\begin{proof}
$N_0$ 为\Term{acyclic-complex}{无上同调复形}，所以\Term{resolution}{分解}为\Term{short-exact-sequence}{短正合列}
\[
0\longrightarrow Z_0^i\longrightarrow N_0^i\longrightarrow Z_0^{i+1}\longrightarrow0,
\]
其中 $Z_0^i$ 为 $N_0^i$ 中的\Term{cycle}{闭元}子\Term{module}{模}；且当 $i$ 充分大时，$N_0^i=0$。对 $i$ 作\Term{descent}{下降}归纳，可见这些\Term{exact-sequence}{正合列}均\Term{split}{分裂}，且每个 $Z_0^i$ 均\Term{projective-module}{投射}。对 $i$ \Term{descent}{下降}归纳，逐步\Term{lifting}{提升}所有这些\Term{exact-sequence}{正合列}，即\Term{liftable}{可提升} $N_0$。
\SourcePage{490}
因此只须证明：若
\[
S_0:\quad0\longrightarrow P_0\longrightarrow Q_0\longrightarrow R_0\longrightarrow0
\]
是 $A_0$-\Term{module}{模}的\Term{exact-sequence}{正合列}，$Q_0$ \Term{projective-module}{投射}，且 $R$ 为满足 $T(R)\simeq R_0$ 的 $A$-\Term{module}{模}，则存在 $A$-\Term{module}{模}的\Term{exact-sequence}{正合列}
\[
S:\quad0\longrightarrow P\longrightarrow Q\longrightarrow R\longrightarrow0,
\]
其中 $Q$ \Term{projective-module}{投射}，且 $T(S)\simeq S_0$。由\Term{lemma}{引理} 3 (ii)，可将 $Q_0$ \Term{lifting}{提升}为\Term{projective-module}{投射} $A$-\Term{module}{模} $Q$；再由\Term{lemma}{引理} 3 (i)，将 $Q_0\to R_0$ \Term{lifting}{提升}为 $A$-\Term{linear}{线性}映射 $Q\to R$。由 \Term{nakayama-lemma}{Nakayama 引理}，该映射\Term{surjective}{满射}；置 $P=\operatorname{Ker}(Q\to R)$，即可补全\Term{exact-sequence}{正合列} $S$。
\end{proof}
\begin{proof}[\Term{lemma}{引理} 2 的证明]
由于 $M=T(L)$ 各次均\Term{projective-module}{投射}，\Term{lemma}{引理}的假设\Term{equivalence}{等价}于说 $T(\varphi)$ 与 $\varphi_0$ \Term{homotopic}{同伦}。因此存在\Term{homotopy-operator}{同伦算子} $k_0=(k_0^i)_i$，$k_0^i:M^i\to M^{i-1}$，使
\[
T(\varphi)-\varphi_0=d_0k_0+k_0d_0,
\]
其中 $d_0$ 为 $M$ 的\Term{differential}{微分}。由\Term{lemma}{引理} 3 (i)，$k_0^i$ \Term{liftable}{可提升}为\Term{linear}{线性}映射 $k^i:L^i\to L^{i-1}$，从而得到 $L$ 中的\Term{homotopy-operator}{同伦算子} $k=(k^i)_i$。只须置 $\varphi'=\varphi-kd-dk$（$d$ 为 $L$ 的\Term{differential}{微分}），便得\Term{lemma}{引理} 2。
\end{proof}
\begin{proof}[\Term{lemma}{引理} 1 的证明]
可以用一个与 $L$ \Term{quasi-isomorphism}{拟同构}的\Term{complex}{复形}替代 $L$，从而假设 $L$，因而也假设 $T(L)$，
\SourcePage{491}
各次均\Term{projective-module}{投射}。下文即作此假设。于是可假设 $D^-(A_0)$ 中的\Term{isomorphism}{同构} $u$ 来自一个仍记作 $u$ 的\Term{quasi-isomorphism}{拟同构} $T(L)\to M$。先\Term{reduced}{约化}到 $u$ \Term{surjective}{满射}的情形。由
\[
N_0^i=M^i\oplus M^{i-1},\qquad
 d^i(x,y)=(d^ix,x-d^{i-1}y)\quad(x\in M^i,\ y\in M^{i-1})
\]
定义的\Term{complex}{复形} $N_0$，显然\Term{acyclic-complex}{无上同调}且各次均\Term{projective-module}{投射}；每次取\Term{projection}{投影} $N_0^i\to M^i$，定义\Term{complex}{复形}的\Term{epimorphism}{满态射} $N_0\to M$。\Term{lemma}{引理} 3 的推论允许将 $N_0$ \Term{lifting}{提升}为右\Term{bounded}{有界}、各次\Term{projective-module}{投射}的\Term{acyclic-complex}{无上同调复形} $N$。于是可将 $L$ 替换为 $L\oplus N$，后者与 $L$ \Term{quasi-isomorphism}{拟同构}（因为 $N$ \Term{acyclic-complex}{无上同调}）且各次均\Term{projective-module}{投射}；同时将 $u$ 替换为\Term{quasi-isomorphism}{拟同构}
\[
T(L\oplus N)=T(L)\oplus N_0\longrightarrow M,
\]
在 $T(L)$ 上取 $u$，在 $N_0$ 上取\Term{epimorphism}{满态射} $N_0\to M$。显然后一个\Term{quasi-isomorphism}{拟同构}是\Term{surjective}{满射}。

因此，设 $u:T(L)\to M$ 为\Term{surjective}{满射}\Term{quasi-isomorphism}{拟同构}，$R_0$ 为其核。由于 $u$ 为\Term{quasi-isomorphism}{拟同构}，$R_0$ \Term{acyclic-complex}{无上同调}；它还各次均\Term{projective-module}{投射}，因为 $M$ 各次均\Term{projective-module}{投射}，故\Term{exact-sequence}{正合列}
\[
0\longrightarrow R_0\longrightarrow T(L)\longrightarrow M\longrightarrow0
\]
是\Term{split}{分裂}的。因此由\Term{lemma}{引理} 3 的推论，可将 $R_0$ \Term{lifting}{提升}为右\Term{bounded}{有界}、\Term{acyclic-complex}{无上同调}且各次均\Term{projective-module}{投射}的 $A$-\Term{module}{模}\Term{complex}{复形} $R$。现证明\Term{morphism}{态射} $f_0:R_0\to T(L)$ \Term{liftable}{可提升}为\Term{complex}{复形}\Term{morphism}{态射} $f:R\to L$。注意 $R_0$ 在\Term{homotopic}{同伦}意义下平凡，所以 $f_0$ 与零\Term{homotopic}{同伦}，可写为
\[
f_0=h_0d_{R_0}+d_{T(L)}h_0,\qquad
h_0=(h_0^i)_i,\quad h_0^i:R_0^i\longrightarrow T(L)^{i-1}.
\]
\Term{lemma}{引理} 3 (i) 允许将 $h_0$ \Term{lifting}{提升}为 $h=(h^i)_i$，其中 $h^i:R^i\to L^{i-1}$。
\SourcePage{492}
可见 $f=hd_R+d_Lh$ 是\Term{lifting}{提升} $f_0$ 的\Term{complex}{复形}\Term{morphism}{态射}。一个标准\Term{lemma}{引理}（参见 SGA 60/61，IV 5.7）表明，$f$ 是\Term{injectivity}{单射}，且 $\operatorname{Coker}f=L'$ 各次均\Term{flat}{平坦}。由于
\[
T(L')\simeq\operatorname{Coker}T(f)=\operatorname{Coker}f_0\simeq M
\]
各次均\Term{projective-module}{投射}，\Term{lemma}{引理} 3 (ii) 给出 $L'$ 也各次均\Term{projective-module}{投射}。另一方面，\Term{canonical}{典范}\Term{homomorphism}{同态} $L\to L'$ 是\Term{quasi-isomorphism}{拟同构}（它\Term{surjective}{满射}且核 $R$ \Term{acyclic-complex}{无上同调}），所以 $L'$ 满足\Term{lemma}{引理} 1 要求的结论。证毕。
\end{proof}
